\documentclass[usenatbib,twoside,twocolumn]{mnras}

% Paper packages
\usepackage{graphicx}
\usepackage{amsmath}
\usepackage{amssymb}
\usepackage{booktabs}

% Denser float packing (reduce float-only pages; no content change)
\renewcommand{\topfraction}{0.9}
\renewcommand{\bottomfraction}{0.85}
\renewcommand{\textfraction}{0.07}
\renewcommand{\floatpagefraction}{0.75}
\renewcommand{\dbltopfraction}{0.9}
\renewcommand{\dblfloatpagefraction}{0.75}
\setcounter{topnumber}{3}
\setcounter{bottomnumber}{2}
\setcounter{totalnumber}{5}
\setcounter{dbltopnumber}{3}

% Define astronomical units
\newcommand{\msun}{\ensuremath{M_{\odot}}}
\newcommand{\pc}{\ensuremath{\rm pc}}
\newcommand{\kms}{\ensuremath{\rm km\,s^{-1}}}

\title[Core spacing and filament occupancy in HGBS filaments]{Local Core Spacing and Filament Occupancy in \textit{Herschel} Gould Belt Filaments\thanks{Results are established for the nearby, low-mass Gould Belt sample ($N=4$ clouds) and are not offered as a universal fragmentation law.}}

\author[G.~J.~White]{G.~J.~White$^{1,2}$\\
$^{1}$School of Physical Sciences, The Open University, Walton Hall, Milton Keynes, MK7 6AA, UK\\
$^{2}$Rutherford Appleton Laboratory, Harwell Science and Innovation Campus, Didcot, OX11 0QX, UK}

\begin{document}

\maketitle

\begin{abstract}
Dense cores along molecular-cloud filaments are often reported at separations shorter than the fastest-growing mode of an equilibrium isothermal cylinder, $\lambda_{\rm m}\approx5$--$6$ filament widths. We re-analyse four \textit{Herschel} Gould Belt Survey clouds with \textit{Gaia} DR3 distances and show that the size of this discrepancy depends on the statistic used.

The median of all pairwise separations is not a local-spacing estimator: it converges on a fixed fraction of the segment length it is applied to. Using adjacent along-crest separations, cores within populated stretches lie only $\sim1$--$3$ widths apart. The filament length per core, the closest analogue of a one-fragment-per-wavelength prediction, is instead $3.4$--$5.8$: one cloud is consistent with the reference band, a second lies just below it, and two are lower by up to a factor of two.

The two diverge because cores are confined to a minority of the crest: on a $\pm0.06$\,pc operational definition only $6$--$27$ per cent is core-bearing, and the inhomogeneity is independently supported by scale-dependent Fano statistics. The gap distribution is inconsistent with a periodic comb, including a deliberately imperfect one, and with a Poisson process, while showing no signature at the classical scale. Core formation is therefore spatially intermittent rather than described by a single wavelength.

Simulations show that a contracting filament need not select the equilibrium wavelength, mode selection tracking instead the density at which the modes grow. They do not reproduce the observed point process, and the one field configuration that could shorten the axial wavelength, a force-balanced helical equilibrium, is not simulated here.
\end{abstract}

\begin{keywords}
stars: formation -- ISM: structure -- ISM: filaments -- MHD -- turbulence
\end{keywords}

\section{Introduction}

The \textit{Herschel} Space Observatory established the ubiquity of filamentary structure in nearby star-forming clouds \citep{Andre2010} and a characteristic width of $\sim$0.1\,pc across diverse environments \citep{Arzoumanian2011}, though the universality and physical interpretation of that scale remain contested \citep{Panopoulou2017,Panopoulou2022}. For an isothermal cylinder the critical mass per unit length is $M_{\rm line,crit}\approx16\,M_\odot\,{\rm pc^{-1}}$ \citep{Ostriker1964}, and a marginally-stable cylinder fragments at $\lambda_{\rm m}\approx5$--$6\times$FWHM \citep{Inutsuka1992,Fischera2012}. Observed core separations are widely reported to fall below that value, and the deficit has been read as evidence for additional physics.

\smallskip\noindent\textbf{What this paper does.} We re-measure core spacing along filament crests in four HGBS clouds with \textit{Gaia} DR3 distances \citep{Zhang2023}, and argue that the apparent conflict with classical theory is largely a question of which statistic is compared with which prediction.

Two observational points anchor the paper. First, the all-pairs (pairwise) median separation is not an estimator of local spacing at all. It converges to $\approx0.293$ times the extent of whatever segment it is applied to, so it measures the region selected (Section~\ref{sec:statistics}). We therefore use adjacent along-crest separations for local spacing and retain the all-pairs statistic only for historical comparison. We make no claim that published HGBS spacings generally suffer from this: the values we could check are local or along-fibre statistics that agree with ours (Table~\ref{tab:literature}).

Second, and more importantly, the local and the global statistics give different answers, and that difference is the result. Measured only within core-bearing stretches, adjacent separations are several times shorter than $\lambda_{\rm m}$ under every skeleton and width convention we tried. Measured over the whole crest, the filament length per core is close to classical in two of the four clouds. The reconciliation is that cores occupy only $6$--$27$ per cent of the crest: they form in tight groups separated by long bare stretches. The scientific question therefore shifts from ``why is the fragmentation wavelength too short?'' to ``why is core formation spatially intermittent along the filament?''.

\smallskip\noindent\textbf{Candidate explanations.} Three mechanisms could shorten a local spacing. (i) \emph{Hierarchical fragmentation}, in which filaments are bundles of velocity-coherent fibres \citep{Hacar2013,Hacar2018,Yang2024}, so that measuring across the bundle compresses the apparent filament-level spacing. (ii) \emph{Magnetic tension} \citep{Nakamura1993}: only a helical or toroidal field shortens the axial mode, since a static longitudinal field leaves it unchanged and a perpendicular field lengthens it. (iii) \emph{Dynamical fragmentation}, in which a still-collapsing filament fragments at a scale set by its current contracting state rather than by the marginally-stable equilibrium mode; this is already qualitatively expected in the dynamical-ISM paradigm \citep{Andre2014,Clarke2016,Clarke2017,Clarke2020}.

The evidence here is not evenly distributed among them. We test (iii) directly, treat (ii) only with a geometric toy model, and do not test the force-balanced helical configuration at all, so (ii) is mentioned for completeness as an unconstrained alternative rather than as a coequal candidate. The data cannot discriminate between the three, and nothing below should be read as identifying the cause.

\smallskip\noindent\textbf{The question this paper asks.} Stated in the form we think the data actually address it: \emph{is core formation along molecular-cloud filaments adequately described by a single fragmentation wavelength, or is it intrinsically spatially intermittent?} The classical-cylinder comparison then becomes one diagnostic of that more general question rather than the organising principle, and $S_{\rm local}$, $S_{\rm global}$, the occupancy fraction, the Fano factor, the gap distribution and the pair-correlation function become different projections of the same underlying point process.

\smallskip\noindent\textbf{New contributions.} These are three: a homogeneous, \textit{Gaia}-DR3, multi-convention measurement of adjacent-core spacing and filament length per core across the four primary clouds, together with the occupancy statistics that connect them; the identification of the region-extent bias in all-pairs separation statistics in this context (the underlying triangular-distribution order statistic is elementary and we claim no new mathematics); and a controlled numerical experiment in which the identical pipeline recovers the classical wavelength for a relaxed equilibrium cylinder but finds no resolved maximum for a contracting one, extended here to a five-point ladder in contraction rate.

\smallskip\noindent\textbf{Scope.} With four independent clouds these are sample statements, not a universal fragmentation law, and all $\pm$ values quoted in this paper are descriptive $N=4$ ranges rather than confidence intervals. Symbols are collected in Table~\ref{tab:symbols}.

\begin{table}
\caption{Principal symbols.}
\label{tab:symbols}
\footnotesize
\setlength{\tabcolsep}{4pt}
\begin{tabular}{lp{6.1cm}}
\toprule
Symbol & Meaning \\
\midrule
$s_{\rm ACS,med}$ & observed adjacent-core spacing (ACS) statistic (median adjacent gap; a line-density scale, \emph{not} a resonant/fragmentation wavelength; cf.\ Section~\ref{sec:clustering}) \\
$W_{\rm fil}$ & observed filament FWHM (beam-deconvolved Plummer; the reference width) \\
$W_{\rm core}$ & simulation initial Gaussian half-width $\sigma$ ($0.3\,\lambda_{\rm J}$) \\
$W_{\rm form}$ & FWHM of the evolved, projected simulation profile \\
$T_1\!\equiv\!\eta_W$ & $W_{\rm form}/W_{\rm fil}$, the simulation-to-observation width normalisation (forward-model band $0.59$--$0.78$; adopted $0.61$) \\
$\lambda_{\rm J}$ & Jeans length at the \emph{initial background} density $\rho_0$ ($=1$ in code units) and never at the evolving central density, for which we write $\lambda_{\rm J}(\rho_c)$ \\
$\lambda_{\rm peak}$ & wavelength at which $\Gamma(\lambda)$ is maximal. Resolved \emph{only} for the relaxed equilibrium control (at $22H$); \emph{not} resolved for any contracting run, for which we quote $\lambda_{\rm turn}$ instead \\
$\lambda_{\rm turn}$ & upper bound on where a turnover in $\Gamma(\lambda)$ could lie for a contracting filament ($\lesssim1.1\,\lambda_{\rm J}$); the only wavelength we quote for those runs \\
$\lambda_{\rm m}$ & equilibrium-cylinder reference wavelength of \citet{Inutsuka1992} (written $\lambda_{\rm IM92}$ where the equilibrium derivation is meant) $\approx5$--$6\,W_{\rm FWHM}$ ($\approx4\times$ flat diameter) \\
$\Gamma(\lambda)$ & linear growth rate of axial mode of wavelength $\lambda$ \\
$f$ & thermal line-mass fraction $M_{\rm line}/M_{\rm line,crit}$ \\
$\beta$ & plasma beta $8\pi P/B^2$ \\
$\mathcal{M}$ & amplitude label of the seed perturbation ($\delta v=\mathcal{M}c_s\times10^{-4}$; a linear seed, \emph{not} a turbulent Mach number; the transonic-perturbation comparison is given in Section~\ref{sec:sims}) \\
$\mu_\Phi/\mu_{\Phi,\rm crit}$ & normalised mass-to-flux ratio \\
$\theta$ & angle between field and filament axis \\
$N_J$ & cells per local Jeans length (Truelove criterion) \\
$C$ & density contrast $\rho_{\rm max}/\langle\rho\rangle$ \\
\bottomrule
\end{tabular}
\end{table}

\label{sec:intro_notation}
\noindent\textbf{Notation.} Four normalised quantities carry the argument. Capital symbols denote ratios to the classical benchmark; $s_{\rm ACS,med}$ denotes a raw separation in parsecs.
\begin{itemize}\setlength{\itemsep}{1pt}
\item $S_{\rm global}=(L/N)/W$, the unconditional filament length per core. This is the closest unconditional observational analogue of the one-fragment-per-wavelength prediction, \emph{provided} the measured crest corresponds to a continuously fragmentation-eligible physical filament. For the HGBS skeletons that proviso does not hold: three lengths must be distinguished, $L_{\rm skel}/N$, $L_{\rm core-bearing}/N$ and $L_{\rm unstable}/N$, and only the last maps directly onto the ideal cylinder. The comparison of $S_{\rm global}$ with $\lambda_{\rm m}/W\approx5$--$6$ is therefore diagnostic rather than a test.
\item $S_{\rm local}=s_{\rm ACS,med}/W$, the median adjacent-core separation measured \emph{only within core-bearing stretches}: a descriptor of the interior of core-bearing regions, not a fragmentation wavelength.
\item $S_{\rm elig}=(L_{\rm elig}/N_{\rm elig})/W$, $S_{\rm global}$ restricted to filament above the critical line mass; an exploratory diagnostic (Section~\ref{sec:eligible}).
\item $S_{\rm sim}^{\rm upper}=\lambda_{\rm turn}/W$, an \emph{upper bound} from the simulations.
\end{itemize}
The difference between $S_{\rm global}$ and $S_{\rm local}$ is itself the measurement of how unevenly cores occupy the crest. Every other ratio quoted here is a sensitivity test or alternative convention on one of these four, decoded in Table~\ref{tab:headline}.

\section{Observational analysis}

\subsection{The sample}

Our sample comprises 8 high-confidence HGBS regions (four \emph{primary}, Table~\ref{tab:sample}; four \emph{limited}, Appendix~\ref{sec:limited_sample}). The original catalogues \citep{Konyves2015, Konyves2020, Ladjelate2020, Pezzuto2021} adopted distances of $130$--$400$\,pc from the best estimates available in 2010--2015 (Table~\ref{tab:literature} lists the pre-Gaia distance and its source per region; note that \citealt{Konyves2020} already adopted $\sim400$\,pc for Orion~B); Gaia DR3 parallaxes \citep{Zhang2023} have since enabled substantial revisions, largest for the more distant regions, whose pre-Gaia distances were the most uncertain.

We adopt the Gaia DR3 YSO-based distances from \citet{Zhang2023} (Table~4), with typical uncertainties of 3\%--5\%. The revisions are strongly region-dependent (largest for Serpens $260\to458$\,pc, $+76\%$, and Aquila $260\to436$\,pc, $+68\%$; Orion~B, Taurus and Ophiuchus essentially unchanged); since spacings scale linearly with distance, the core-count-weighted mean rises only modestly ($\sim10\%$ for the primary regions), dominated by Aquila rather than being a uniform $+32\%$ shift. Because Gaia DR3 thus compresses or expands each cloud non-uniformly, the revision is not a single global rescaling, which is a further reason the clouds must be treated as $N=4$ independent units rather than by pooling cores that carry very different distance corrections. The observed spacing is a line-density statistic, not a resonance, and coincides with the theoretical wavelength only to order of magnitude.

The analysis includes all HGBS-identified cores (starless, prestellar and protostellar), extracted and classified within the consistent HGBS framework \citep{Konyves2015, Konyves2020}. Including unbound starless and protostellar sources may introduce some bias, but fragmentation produces a continuum of boundness states, and restricting the sample to prestellar cores would shrink several regions (notably TMC1, CRA and Serpens) below statistical usefulness. Core positions and properties are taken directly from the published catalogues.

Protostellar migration displacements of $0.01$--$0.05$\,pc \citep{Kirk2016, Mattern2018} are small against the $\sim0.2$--$0.3$\,pc spacing; we estimate the resulting ACS bias at $\sim5$--$10\%$ and carry it as a systematic (Section~\ref{sec:error_budget}).

\begin{table}
\caption{Primary HGBS sample with Gaia DR3 distances. Per-region spacings are pairwise-median values rounded to three significant figures; the core-count-weighted mean of the four primary regions is $0.290$\,pc. The four \emph{limited} regions (Ophiuchus, Serpens, TMC1, CRA) and the full 8-region weighted mean ($0.279$\,pc) are listed in Appendix~\ref{sec:limited_sample}.}
\label{tab:sample}
\footnotesize
\setlength{\tabcolsep}{3pt}
\begin{tabular}{lccccc}
\toprule
Region & $D$ & $N_{\rm core}$ & Pairwise med. & S.E. & Status \\
       & (pc) &  & (pc) & (pc) &  \\
\midrule
Orion B  & 386 & 1,844 & 0.313 & 0.047 & Primary \\
Aquila   & 436 & 749   & 0.346 & 0.047 & Primary \\
Perseus  & 300 & 816   & 0.248 & 0.040 & Primary \\
Taurus   & 135 & 536   & 0.198 & 0.040 & Primary \\
\midrule
\textbf{Weighted Mean}$^\dagger$ &  & \textbf{3,945} & \textbf{0.290} &, &  \\
\bottomrule
\end{tabular}
\par\vspace{2pt}
\footnotesize $^\dagger$Core-count-weighted mean of the four primary regions. A formal core-count-weighted standard error (treating cores as independent) is \emph{not} the relevant uncertainty, since the four clouds are the independent units (Section~\ref{sec:error_budget}); the cloud-to-cloud scatter is carried instead.
\end{table}

\subsection{Measurement method}
\label{sec:methodology}

The analysis proceeds from the HGBS column-density maps and \textit{Gaia} DR3 distances through skeleton extraction, morphological bridging, core association within $0.06$\,pc of a spine, graph ordering along each component and adjacent-core separation. Filament skeletons are the published HGBS DisPerSE products \citep{Sousbie2011,Arzoumanian2011}, extracted with a $3\sigma$ persistence threshold and manual editing restricted to removing obvious artefacts. The effect of that editing on the spacing is small ($\lesssim5$ per cent), bounded by the junction-retention and association-threshold sensitivity tests; robustness to the persistence threshold itself is examined in Appendix~\ref{sec:persistence}. Because skeletons and core positions derive from distance-independent angular positions, the \textit{Gaia} revisions require no re-extraction. Excluding junctions changes the measured spacing by $<5$ per cent and varying the association threshold between one and $1.5$ widths by $<3$ per cent.

Core-to-core spacings were also measured with the pairwise median, retained only as a comparability statistic with earlier work (Section~\ref{sec:statistics}).

\smallskip\noindent\textbf{Distances.} Physical spacing is angular separation times distance, so the $\lambda\propto d$ scaling is tautological and we do not present it as corroboration. The relevant check is that the residuals about that scaling show no trend with the absolute revision fraction, and the two most-revised clouds do not stand out. Two caveats remain. The YSO-based distances \citep{Zhang2023} are assumed to apply to the dense gas measured here; and the Aquila/Serpens Rift has substantial line-of-sight depth \citep{Zucker2020,Green2024}, so a single assigned distance may blend populations.

We quantified the blending with a geometric Monte-Carlo model in which two crests at different depths are projected onto a common spine, taking an intrinsic spacing of $0.20$\,pc at $d=436$\,pc. Each crest's offset is drawn from a Gaussian of dispersion $\sigma_z$, with the two separated by $2\sigma_z$, scanning $\sigma_z=5$--$30$\,pc. Superposition interleaves the two chains and \emph{shortens} the measured median, to $0.41$--$0.49$ of the true value across that range, whereas the distance spread alone contributes only $\pm1$ to $\pm7$ per cent. Blending therefore biases the spacing towards shorter values rather than longer. The $\sim50$ per cent figure is a worst case in which two crests overlap completely. At the plausible end, $\sigma_z=5$--$10$\,pc with half the crest length overlapping, the shortening factor is $0.78$--$0.86$, which would move Aquila's $S_{\rm local}$ from $1.7$ to $2.0$--$2.2$ and $S_{\rm global}$ from $3.4$ to $4.0$--$4.4$: still sub-classical locally and still below the reference band globally, so Aquila's classification survives a realistic correction. This is a geometric model with a free depth parameter, not a measurement; the stronger treatment, projecting the three-dimensional dust reconstructions along the actual sightline, is listed as a follow-up in Section~\ref{sec:limitations}. As a floor, retaining the pre-\textit{Gaia} distances lowers the primary-region characteristic to $S_{\rm local}\approx1.4$, so the sign of the result does not depend on the distance revision though its magnitude does.

\begin{figure}
\centering
\includegraphics[width=0.99\columnwidth]{figures/fig_skeleton_compare.pdf}
\caption{The dominant observational systematic, shown directly: the DisPerSE (blue) and FilFinder (red) skeletons extracted from the \emph{same} \textit{Herschel} column-density map of Orion~B (greyscale, $\log N_{\rm H_2}$; a representative sub-region). The FilFinder skeleton is produced by \textsc{FilFinder2D} with the cloud distance set per region and an \emph{externally supplied} mask (\texttt{use\_existing\_mask=True}) followed by \texttt{medskel()}; the mask is a simple column-density cut at the $92$nd percentile of the finite pixels, which corresponds to $N_{\rm H_2}=2.5$, $3.8$, $3.1$ and $2.9\times10^{21}$\,cm$^{-2}$ in Orion~B, Aquila, Perseus and Taurus respectively. Because the mask is supplied externally, \textsc{FilFinder}'s own adaptive and branch thresholds are not invoked. A more permissive $70$th-percentile mask is also computed and is what sets the upper end of the FilFinder range. The DisPerSE skeletons are the published HGBS products, extracted at a $3\sigma$ persistence cut. The two algorithms trace substantially different filament networks, FilFinder's medial axis of a column-density mask is denser and more branched and follows more low-contrast material, whereas DisPerSE follows the persistent crests, which is why the inferred $s_{\rm ACS,med}/W$ differs by a factor of $\sim$2 between them (Section~\ref{sec:error_budget}). This is intrinsic to how each algorithm defines a ``filament'' and is the largest single uncertainty on the \emph{magnitude} of the spacing, though the sub-classical \emph{sign} holds for both.}
\label{fig:skeleton_compare}
\end{figure}

\subsection{Results}
\label{sec:results}

The four primary regions are the independent statistical units throughout; the nesting of cores within clouds is propagated in the hierarchical bootstrap of Section~\ref{sec:error_budget}.

\smallskip\noindent\textbf{Two statistics, two different answers.} The adjacent-core spacing $s_{\rm ACS,med}$ is \emph{conditional}: measured only where cores are adjacent, it characterises the interior of core-bearing stretches and says nothing about filament that has formed no cores. The core line density $N/L$, quoted through its reciprocal $L/N$, is \emph{unconditional}.

Taking the unconditional quantity first, $S_{\rm global}=3.4$, $3.4$, $4.5$ and $5.8$ for Orion~B, Aquila, Taurus and Perseus. Against the equilibrium-cylinder reference band, Perseus is essentially consistent with it, Taurus lies just below, and Orion~B and Aquila are below it by less than a factor of two. Conditionally, $S_{\rm local}$ is shorter by a factor $\approx2.4$--$4.3$ (Table~\ref{tab:headline}). The reconciliation is the occupancy result of Section~\ref{sec:intermittency}.

Every comparison of $S_{\rm local}$ with the reference band below is therefore between a conditional descriptor and an equilibrium wavelength, and is diagnostic rather than a test. The sign of the local deficit is robust under every convention we tried; its magnitude is not, being dominated by the choice of skeleton. Figure~\ref{fig:envelope} shows the full methodological envelope on a single axis and is the quickest orientation to the sections that follow.

\begin{table}
\caption{Principal results, one row per cloud. $S_{\rm local}=s_{\rm ACS,med}/W$ is the conditional adjacent-core spacing, measured only within core-bearing stretches; $S_{\rm global}=(L/N)/W$ is the unconditional filament length per core; $S_{\rm elig}=(L_{\rm elig}/N_{\rm elig})/W$ restricts the latter to filament above the critical line mass, selected on a core-masked column-density reconstruction (Section~\ref{sec:eligible}), with $f_{\rm elig}$ the eligible length fraction and $N_{\rm elig}$ the cores on it; $f_{\rm occ}$ is the $0.06$\,pc occupancy fraction. All are normalised by the fixed $W=0.10$\,pc width so the columns are directly comparable, and all use the DisPerSE skeleton. The classical benchmark is $5$--$6$. The simulated upper bound is common to all rows, $S_{\rm sim}^{\rm upper}=\lambda_{\rm turn}/W\lesssim1.0$--$1.3$ (Section~\ref{sec:sims}). $^{\dagger}$~The $S_{\rm elig}$ and $f_{\rm elig}$ columns are an \emph{exploratory diagnostic}, not a measurement: the eligibility criterion is a column-density proxy for a line mass and is partly determined by the cores it is used to interpret (Section~\ref{sec:eligible}); they should not be read on the same footing as $S_{\rm local}$ and $S_{\rm global}$. The Perseus entry is bracketed because only $57$ cores remain on eligible filament after masking and the value depends on the masking radius (Table~\ref{tab:maskrad}); no value is quoted for Taurus, where $13$ cores remain and the statistic does not stabilise.}
\label{tab:principal}
\footnotesize
\setlength{\tabcolsep}{4pt}
\begin{tabular}{lccccc}
\toprule
Cloud & $S_{\rm local}$ & $S_{\rm global}$ & \emph{$S_{\rm elig}$}$^{\dagger}$ & \emph{$f_{\rm elig}$ [$N_{\rm elig}$]}$^{\dagger}$ & $f_{\rm occ}$ \\
\midrule
Orion~B & 1.8 & 3.4 & \emph{2.0}$^{\dagger}$ & \emph{0.16 [268]} & 0.27 \\
Aquila  & 1.7 & 3.4 & \emph{3.2}$^{\dagger}$ & \emph{0.90 [303]} & 0.24 \\
Perseus & 3.0 & 5.8 & \emph{(4.1)}$^{\dagger}$ & \emph{0.08 [57]} & 0.11 \\
Taurus  & 1.8 & 4.5 & ---    & --- [13]   & 0.06 \\
\bottomrule
\end{tabular}
\end{table}

\begin{table}
\caption{Convention decoder for the \emph{conditional} local descriptor (symbols: Table~\ref{tab:symbols} and the notation paragraph of Section~\ref{sec:intro_notation}): read every observed $s_{\rm ACS,med}/W$ quoted in the text against it. This table tabulates the adjacent-core spacing, which is measured only within core-bearing stretches and is \emph{not} the quantity a one-fragment-per-wavelength theory predicts (that is $L/N$; Section~\ref{sec:linedensity}). Values are given under each width/skeleton convention and correction state for the four primary HGBS regions (revised distances); every combination is conditionally sub-classical (classical $\lambda_{\rm m}\approx5$--$6\times$FWHM). We adopt the directly-measured region-Plummer value as our reference, $\boldsymbol{(s_{\rm ACS,med}/W)_0\approx1.4}$; the fixed-width value ($\approx1.9$, retained for comparability with prior HGBS work) and all others are read against it. The ``periodic'' column applies the periodic injection--recovery correction to the raw value; the data-matched \emph{clustered} column (bracketed) applies $\div1.64$. The DisPerSE conventions are sub-classical by a factor $\approx2.4$--$4.3$; across the full envelope the factor spans $\approx2$--$13$, the upper end set by the densest FilFinder skeleton.}
\label{tab:headline}
\footnotesize
\setlength{\tabcolsep}{4pt}
\begin{tabular}{p{4.6cm}cc}
\toprule
Quantity & raw & periodic \\
\midrule
\multicolumn{3}{@{}l}{\emph{Conditional local descriptor} $s_{\rm ACS,med}/W$ (core-bearing stretches only)}\\
Region Plummer, DisPerSE (adopted, $(s_{\rm ACS,med}/W)_0$) & $\mathbf{1.4}$ & $1.2$ \\
Fixed $0.10$\,pc, DisPerSE & $2.1$ & $1.9$ \\
FilFinder skeleton & $0.4$--$0.8$ & $0.4$--$0.8$ \\
Per-segment local width, $\Delta s/W_{\rm local}$ & \multicolumn{2}{c}{$0.6$--$1.1$} \\
\emph{Deprojected} (adopted row $\div\,\langle\sin i\rangle_{\rm med}=0.87$) & $1.6$ & $1.4$ \\
\emph{Deprojected} (fixed $0.10$\,pc) & $2.4$ & $2.2$ \\
\midrule
\multicolumn{3}{@{}l}{\emph{Unconditional} (what fragmentation theory predicts)}\\
Global length per core $(L/N)/W$, fixed $0.10$\,pc & \multicolumn{2}{c}{$3.4$--$5.8$} \\
Mean adjacent, fixed $0.10$\,pc (branch-jump inflated) & \multicolumn{2}{c}{$2.1$--$4.7$} \\
\midrule
\multicolumn{3}{@{}l}{\emph{Legacy comparator and reference scales}}\\
Pairwise median ($\approx0.293L$; biased, not a spacing) & $2.8$ & \\
Matched simulation (upper limit) & \multicolumn{2}{c}{$\lambda_{\rm turn}/W\lesssim1.0$--$1.3$} \\
Classical $\lambda_{\rm m}$ (F\&M 2012) & \multicolumn{2}{c}{$5$--$6$ ($=4\times$ diam.)} \\
\bottomrule
\end{tabular}
\\[2pt]
\footnotesize
\emph{Notes} The reference point is the region-Plummer ACS $(s_{\rm ACS,med}/W)_0\approx1.4$ (bold), one illustrative point within the convention range $0.8$--$1.9$; it is an adopted reference value, chosen for comparability with prior HGBS work, not a uniquely preferred physical scale (and not a formal confidence interval; Section~\ref{sec:error_budget}). The data-matched \emph{clustered} branch ($\div1.64$ on the raw value, giving $\approx0.8$ region-Plummer, $\approx1.2$ fixed-width and $\approx0.3$--$0.5$ FilFinder) is \emph{not} shown as a column because it is not independent: it applies a correction derived by \emph{assuming} the observed clustering (Section~\ref{sec:injection_recovery}), so it is circular and serves only as a lower bound; we carry the raw value throughout. The fixed-width convention gives a \emph{higher} value ($\approx1.9$) because it uses $0.10$\,pc rather than the larger measured widths; FilFinder gives lower values still. The full envelope factor $\approx2$--$13$ (Section~\ref{sec:injection_recovery}) is $\lambda_{\rm m}/(s_{\rm ACS,med}/W)$: its lower endpoint $\approx2$ uses $(\lambda_{\rm m},s_{\rm ACS,med}/W)=(5.5W,2.6)$ (fixed-width DisPerSE, Perseus) and its upper endpoint $\approx13$ uses $(5.5W,0.42)$ (densest FilFinder skeleton). We quote $s_{\rm ACS,med}/W$ to one decimal place; finer digits are not significant. The two \emph{deprojected} rows apply the median projection factor $\langle\sin i\rangle_{\rm med}=0.87$ for a \emph{random} distribution of filament orientations. That assumption need not hold within a single cloud complex, where a coherent magnetic field can align the filaments, so the correction is indicative rather than exact (Section~\ref{sec:error_budget}); this correction acts \emph{against} our conclusion, raising the ratio by $15$--$27$ per cent, and we show it in the table rather than only in the text so that the effect is visible in the same place as the adopted numbers. We nevertheless carry the uncorrected value as the headline, because deprojection assumes an isotropic orientation distribution that the filaments in a single cloud need not obey.
\end{table}

\subsection{The dominant observational systematic: filament width}
\label{sec:width}

Because $S_{\rm local}\propto1/W$, the width convention sets the normalisation. Two recur: the \emph{fixed-width} convention ($W_{\rm fil}=0.10$\,pc for every region) and the \emph{region-Plummer} convention, a beam-deconvolved Plummer FWHM fitted per region. The broader measured widths give a lower ratio, so the fixed-width value is the more conservative of the two.

We fit Plummer profiles perpendicular to the crest at every skeleton pixel, over a half-range of $0.4$\,pc with the outer $0.1$\,pc used to set a local background, and deconvolve the $18.2''$ beam. Every region returns an index $p\simeq2.0$, the value found for HGBS filaments generally \citep{Arzoumanian2019}, and a width of $0.11$--$0.16$\,pc, within $\sim50$ per cent of the canonical $0.10$\,pc and consistently above it (Table~\ref{tab:widths}). The region-specific widths are the physically preferable normalisation and we adopt them as the reference; $0.10$\,pc is retained for comparability with earlier work.

\begin{table}
\caption{Region-specific, beam-deconvolved filament widths from a direct Plummer fit to perpendicular column-density profiles along the skeletons (the HGBS width definition; \citealt{Arzoumanian2019}), and the resulting ACS $s_{\rm ACS,med}/W$. $W_{\rm fil}$ is the beam-deconvolved Plummer FWHM ($18.2''$ beam), fitted over a perpendicular half-range of $0.4$\,pc either side of the crest with the outer $0.1$\,pc used to set a local background, $p$ the Plummer index (a \emph{free} fit parameter, not fixed; it converges to $p\simeq2.00$ in all four regions (Taurus $2.01$, a difference in the second decimal place only), consistent with the canonical Plummer-2 profile), $N$ the number of fitted profiles. Quoted $W_{\rm fil}$ uncertainties are the standard error on the median of the per-profile fits within each region and do not include the beam-deconvolution or background-subtraction systematics, which are carried separately in the error budget (Section~\ref{sec:error_budget}); the fitted $p$ has a per-region dispersion of $\pm0.1$ across profiles. All four regions give $p\simeq2$ and sub-classical $s_{\rm ACS,med}/W$ on the (primary-sample) \emph{median} spacing quoted here; the \emph{mean} adjacent spacing is a factor $\approx1.3$--$2.2$ larger (Section~\ref{sec:clustering}), so the mean-based ratio is $\approx2.1$--$4.7$, still sub-classical in every region.}
\label{tab:widths}
\footnotesize
\setlength{\tabcolsep}{3pt}
\begin{tabular}{p{1.7cm}ccccc}
\toprule
Region & $W_{\rm fil}$ (pc) & $p$ & $N$ & ACS (pc) & $s_{\rm ACS,med}/W$ \\
\midrule
Taurus  & $0.110$ & $2.01$ & 295 & $0.124$ & $1.1$ \\
Orion B & $0.148$ & $2.00$ & 356 & $0.207$ & $1.4$ \\
Perseus & $0.147$ & $2.00$ & 235 & $0.263$ & $1.8$ \\
Aquila  & $0.161$ & $2.00$ & 446 & $0.210$ & $1.3$ \\
\midrule
median  & $0.148$ & $2.00$ &, &, & $1.4$ \\
\bottomrule
\end{tabular}
\par\vspace{2pt}
\footnotesize The measured widths ($0.11$--$0.16$\,pc) agree with the published HGBS values and exceed the canonical $0.10$\,pc, so region-specific Plummer widths give a \emph{lower} $s_{\rm ACS,med}/W$ ($\approx1.4$, raw) than the raw fixed-width value ($\approx2.1$; Table~\ref{tab:headline}): all four regions are sub-classical under either convention. ACS spacings are the full-coverage medians (Section~\ref{sec:injection_recovery}), measured on the same DisPerSE skeleton as the widths (its deposited crest for the width fit, its bridged form for the spacing); the injection--recovery bias correction lowers each $s_{\rm ACS,med}/W$ by a further $\sim$10--20\%. (Pipeline outputs retained to 2 s.f.; the convention systematic dominates.)
\end{table}

\subsection{Sample classification and distance robustness}
\label{sec:distance_uncertainties}

We classify the eight regions (Table~\ref{tab:sample}; Appendix~\ref{sec:limited_sample}) by three criteria applied before any spacing was measured, namely a \textit{Gaia} DR3 distance uncertainty below $10$ per cent, at least $100$ catalogue cores, and a skeleton whose largest connected component carries at least a quarter of the total length, as \textbf{primary} (Orion~B, Aquila, Perseus, Taurus) or \textbf{limited} by three criteria fixed \emph{a priori} (before the measurement): $N>500$ cores; a well-established distance; and a skeleton simple enough for unambiguous core-to-filament association. This selection is cross-checked by the leave-one-out test below. The limited sample fails at least one criterion (Ophiuchus on skeleton complexity, the others on small samples) and are excluded from the primary ACS measurement. The largest distance revisions (Serpens $+76\%$, Aquila $+68\%$) are independently corroborated by VLBI maser parallaxes (validating Gaia DR3 to $<5\%$; \citealt{Kounkel2018,Xu2021,OrtizLeon2018,OrtizLeon2017}), extinction estimates \citep{Yan2022,Green2024} and co-spatiality with Aquila \citep{Zucker2020}; Serpens is \emph{limited} and excluded in any case. The sub-classical conclusion is independent of these regions: leave-one-out shifts the weighted mean by $<6\%$, and dropping Aquila leaves the region-Plummer median at $1.40$ and the fixed-width characteristic at $2.1$, both sub-classical. The primary/limited split is likewise not tuned to the distance revision: its criteria act on the distance-independent core count and skeleton morphology, the closest any region lies to the $N>500$ boundary is Taurus at $536$ cores ($\sim7\%$ above), and because the Gaia revision rescales spacings but not core counts, no plausible distance change moves a region across the classification.

\subsection{Choice of spacing statistic}
\label{sec:statistics}

We adopt the along-skeleton adjacent-core spacing (ACS), denoted $s_{\rm ACS,med}$ for its per-region median, as the primary local measure, and report the all-pairs median only for comparability with prior work. It is not the Euclidean nearest-neighbour distance, which we use only for the Clark--Evans statistic of Section~\ref{sec:clustering}.

An all-pairs median, used as a local-spacing estimator, is intrinsically dominated by the extent of the selected segment. For $N$ points distributed over a segment of length $L$ it converges to $d^\star=L(1-1/\sqrt{2})\approx0.293\,L$, the median of the triangular pair-distance law, for a random and a perfectly periodic distribution alike (Appendix~\ref{sec:pwmedian}). It tracks that value for any filament with $N\gtrsim5$ cores and never approaches the local spacing, which is $30$--$100$ times smaller. The order statistic is elementary; what matters here is its consequence for a statistic in common use.

Synthetic point processes confirm this directly. For cores placed by periodic, clustered (gamma-renewal), hierarchical (three projected fibres) and fractal processes with a fixed injected local scale, the ACS median is flat with segment extent and recovers each process's local scale, whereas the pairwise median rises as $0.29\,L$ for \emph{every} process, independent of the injected scale.

We reconstruct the along-skeleton core ordering for the four primary regions from the deposited DisPerSE skeleton maps, bridging fragmented segments by morphological closing (Section~\ref{sec:methodology}). Figure~\ref{fig:spacing} shows the per-region spacings.

\begin{figure*}
\includegraphics[width=0.95\columnwidth]{figures/figure1_spacing_comparison.pdf}
\caption{The conditional local descriptor $S_{\rm local}=s_{\rm ACS,med}/W_{\rm fil}$ per primary region. This is \emph{not} the quantity classical theory predicts (that is $S_{\rm global}$; Fig.~\ref{fig:intermittency}b) and the comparison with the classical band below should be read diagnostically. It is the \emph{conditional} descriptor measured within core-bearing stretches; the unconditional length per core $L/N$, whose reciprocal is the line density predicted by fragmentation theory, is shown in Fig.~\ref{fig:intermittency}, in the two DisPerSE width conventions: directly-measured region-Plummer widths (blue circles, preferred; per-region raw $1.1,1.4,1.8,1.3$, dashed line at the median, which we quote as $1.4$, the rounded value of the exact median of $1.35$) and the fixed $0.10$\,pc width (orange squares; $1.2,2.1,2.6,2.1$), both shown \emph{raw} (before the injection--recovery correction; Section~\ref{sec:injection_recovery}). The red band is the equilibrium-cylinder reference wavelength $\lambda_{\rm m}\approx5$--$6\times$FWHM, exact for the $p=4$ Ostriker equilibrium; the conversion to FWHM is not invariant for the observed $p\simeq2$ profiles (\citealt{Fischera2012}; $=4\times$ the flat diameter, \citealt{Inutsuka1992}) and the dotted line the classical \emph{critical} wavelength $\lambda_{\rm cr}\approx2.6\times$FWHM; every region lies well below both. The grey band marks the FilFinder skeleton range ($\approx0.4$--$0.8$), and the vertical grey bars give the per-region min--max convention envelope, spanning the width conventions and the skeleton-algorithm factor. These are descriptive envelopes rather than confidence intervals (Section~\ref{sec:error_budget}).}
\label{fig:spacing}
\end{figure*}

\subsection{Injection--recovery validation}
\label{sec:injection_recovery}

To test whether the estimator recovers a known spacing under realistic conditions, synthetic cores with periodic spacings of $0.15$--$0.35$\,pc were placed on HGBS-property-matched skeletons and passed through the identical pipeline. The recovery is biased low by $11$--$29$ per cent, because bridging and association occasionally split or merge a pair, and the correction is applied as a per-region multiplicative factor.

The size of that correction depends on the assumed point process, which matters because the real cores are not periodic. Repeating the injection with a clustered (gamma-renewal) process matched to the observed coefficient of variation gives a larger correction ($\div1.64$). That branch is not independent, since it derives the correction by assuming the clustering it is meant to test, so we carry the periodic-injection value throughout and quote the clustered branch only as a lower bound (Table~\ref{tab:headline}). Sweeping the association half-width over $0.04$--$0.08$\,pc changes the per-region median by $4$--$7$ per cent.

\smallskip\noindent\textbf{Full-coverage measurement.} The measurement uses the dense deposited DisPerSE skeletons for all four primary regions. As a segmentation sensitivity test rather than an independent confirmation, since it uses the same maps and catalogue, we repeated the whole analysis on FilFinder skeletons (Section~\ref{sec:ontology}).



\subsection{Core line density and the local-width test}
\label{sec:linedensity}

The median adjacent spacing measures the separation \emph{where cores are adjacent}. For a filament unstable along its whole measured length and producing one core per unstable wavelength, the corresponding statistic is the core line density $N/L$, quoted through its reciprocal $L/N$. That identification holds only under those conditions: the measured skeleton length is not in general the length of unstable cylinder to which the prediction applies, and a filament half of whose length is subcritical returns $L/N\simeq2\lambda$ even if the unstable half fragments perfectly at $\lambda$. Segmentation therefore changes $L/N$ without changing the physics, which is why $S_{\rm global}$ moves from $3.4$--$5.8$ to $9.1$--$16.7$ between the two skeleton definitions (Section~\ref{sec:ontology}), and why the eligible-length diagnostic of Section~\ref{sec:eligible} follows immediately below.

We report three statistics, all normalised by the same fixed $0.10$\,pc width. (i)~The median adjacent spacing is $S_{\rm local}\approx1.7$--$3.0$, sub-classical in every cloud. (ii)~The mean adjacent spacing is a factor $1.3$--$2.2$ larger, but on heavily bridged skeletons it is inflated by arclength jumps across branches, so we treat it as an upper descriptor. (iii)~The length per core is $L/N=0.34$--$0.58$\,pc, that is $S_{\rm global}\approx3.4$ (Orion~B), $3.4$ (Aquila), $4.5$ (Taurus) and $5.8$ (Perseus), a factor $\approx2$ larger than the median adjacent spacing in every cloud.

Two contributions produce that factor of two, and neither is small-scale clustering. The median-to-mean rise ($\approx40$--$120$ per cent) is the generic skew of a broad gap distribution, a pure exponential already giving $1.44$; the further rise to $L/N$ is dilution by extensive core-free filament. The line density restricted to core-bearing stretches is intermediate and still mildly sub-classical, so the near-classical $S_{\rm global}$ is not evidence of a uniform classical spacing.

A \emph{local-width} normalisation supports the same reading. Fitting the beam-deconvolved transverse width at the midpoint of every adjacent-core interval ($146$--$545$ reliable fits per cloud) gives median $\Delta s/W_{\rm local}=0.61$, $0.91$, $1.05$ and $0.99$ (Taurus, Orion~B, Perseus, Aquila), all sub-classical, with only a weak spacing--width correlation ($r=-0.06$ to $0.31$). These median-of-ratios values are if anything more sub-classical than the corresponding ratio-of-medians, so the per-segment normalisation strengthens the conclusion.



\subsection{An exploratory diagnostic: the column-density-selected eligible length}
\label{sec:eligible}

Identifying $L/N$ with the quantity a one-fragment-per-wavelength theory predicts assumes that the whole skeleton could have fragmented. Much of it could not: a subcritical stretch has no axial instability to express. The estimate below is exploratory: the criterion is a column-density proxy for a line mass, and it is partly determined by the cores it is used to interpret. None of the conclusions of this paper rests on it.

We define eligibility \emph{independently of where the cores are}, using column density. A filament of width $0.10$\,pc reaches the critical line mass $M_{\rm line,crit}=2c_s^2/G\simeq16\,M_\odot\,{\rm pc^{-1}}$ at $N_{\rm H_2}\simeq7.2\times10^{21}\,{\rm cm^{-2}}$ ($\mu=2.8$), and we take the \emph{column-density-selected} eligible length $L_{\rm elig}$ to be the skeleton length above that threshold. This is a proxy, not a line-mass measurement: the conversion depends on width, inclination, temperature, background subtraction and both non-thermal and magnetic support. The eligible fractions differ greatly between clouds, $0.19$, $0.91$, $0.13$ and $0.07$ for Orion~B, Aquila, Perseus and Taurus, so Orion~B in particular is dominated by filament that could not have fragmented at all.

\smallskip\noindent\textbf{The criterion is partly endogenous.} A core raises the column density of its own surroundings, so a threshold applied to the raw map can be crossed \emph{by the core it is meant to be independent of}, selecting short eligible patches centred on cores. Between $25$ and $34$ per cent of skeleton pixels lie within one association half-width of a core, so this is not a small effect. We therefore rebuilt the criterion on a core-masked background: pixels within a masking radius of an associated core have their column density replaced by interpolation along the spine from untainted pixels, and the threshold is applied to that reconstruction.

The correction is substantial. On the core-masked background $S_{\rm elig}$ becomes $2.0$, $3.2$, $4.1$ and $3.1$, against the raw $1.6$, $3.2$, $2.0$ and $0.8$. Aquila is unchanged because almost all of its skeleton is eligible either way, but Perseus doubles and Taurus rises by nearly a factor of four, and $N_{\rm elig}$ falls from $199$ to $57$ in Perseus and from $186$ to $13$ in Taurus.

The statistic is also sensitive to the masking radius itself (Appendix~\ref{sec:maskapp}). Only Orion~B and Aquila are comparatively stable against it and retain more than $230$ eligible cores; Perseus drifts over $3.6$--$4.5$ with a bootstrap uncertainty of $^{+1.1}_{-0.9}$ on $57$ cores, and Taurus does not stabilise at any radius.



Sensitivity to the threshold itself is milder and runs the other way: raising it from $4\times10^{21}$ to $1.2\times10^{22}\,{\rm cm^{-2}}$ moves $S_{\rm elig}$ monotonically downwards, so the fiducial choice is conservative, and adopting each region's own Plummer width or $T=15$\,K changes it by less than $20$ per cent. A virial threshold, $M_{\rm line,vir}=2(c_s^2+\sigma_{\rm nt}^2)/G$, lowers it further, but treating $\sigma_{\rm nt}$ as isotropic support is a prescription rather than a measurement and beyond the transonic case the statistic degenerates.

\smallskip\noindent\textbf{What distinguishes a fertile stretch from a sterile one?} Supercriticality is necessary but not sufficient: in Aquila $90$ per cent of the crest is eligible yet only $24$ per cent carries a core. Local compression, longitudinal velocity gradients of $\sim1$\,km\,s$^{-1}$\,pc$^{-1}$ \citep{Kirk2013,Hacar2013}, and spatially varying non-thermal support are all plausible controls, and none can be tested with column density alone. Deciding among them requires resolved kinematics per stretch, so the occupancy is reported here as a measured property whose cause is undetermined.

The message is limited but useful: a physically motivated restriction to nominally unstable crest moves $S_{\rm local}$ and $S_{\rm global}$ towards each other, but current column-density data do not permit a sufficiently independent determination of the unstable filament length.

\smallskip\noindent\textbf{A prescription for future use.} If this diagnostic is to be comparable between studies the masking radius must be tied to a physical scale rather than chosen. We recommend setting it to one filament FWHM measured in the same region, subject to a floor of one beam, so that $r_{\rm mask}=\max(W_{\rm fil},\theta_{\rm beam}d)$; for our sample that is $0.11$--$0.16$\,pc, at the upper end of the range scanned here, and it is the choice that removes the core's own contribution most completely without depending on the association convention. Reporting the statistic at that radius together with $N_{\rm elig}$, and declining to quote it where $N_{\rm elig}\lesssim100$, would make the diagnostic usable. Separating fertile from sterile supercritical stretches, however, is not a column-density problem at all: it requires dense-gas kinematics, specifically N$_2$H$^+$(1--0) or NH$_3$(1,1)/(2,2) mapping at $\lesssim0.05$\,pc resolution, from which the local non-thermal linewidth gives the virial line mass directly and the velocity gradient along the crest distinguishes stretches that are being assembled from those that are being sheared. Resolved kinematics per stretch would be needed to convert this diagnostic into a measurement, and none of the conclusions below depends on it.

\subsection{Does the local/global contrast survive changing the definition of a filament?}
\label{sec:ontology}

A skeleton defines $L$ itself, so changing the extraction algorithm alters the spacing, the length, the core association and the topology together. DisPerSE and FilFinder are not two noisy estimates of one observable but partly different definitions of the object, and we therefore recomputed \emph{all} the principal statistics under both (Table~\ref{tab:ontology}).

The FilFinder networks are shorter ($67$--$121$\,pc against $106$--$337$\,pc) and associate far fewer catalogue cores, so both statistics move, but they move in \emph{opposite} directions: $S_{\rm local}$ becomes shorter still ($0.54$--$1.54$ against $1.66$--$2.97$) while $S_{\rm global}$ becomes much larger ($9.1$--$16.7$ against $3.4$--$5.8$). The local/global contrast therefore survives the change of ontology and in fact widens, from a factor $\approx2$ to $\approx6$--$17$, and the occupancy fraction stays low under both. We regard this as one of the more secure results here, precisely because the two algorithms disagree so strongly about what a filament is.

At a high-column hub, FilFinder recovers the main ridge but adds long, almost straight spokes that leave the emission entirely (Fig.~\ref{fig:skeleton_compare}). These are most likely medial-axis artefacts of a large connected mask rather than resolved sub-filaments, though we have not quantified what fraction of the total length behaves this way and column density alone cannot decide: position--position--velocity data are needed to establish which network corresponds to a coherent physical object. We therefore base our quantitative conclusions on the persistent-crest network and carry FilFinder as an ontology test.

\smallskip\noindent\textbf{Coverage bias.} Between $41$ and $86$ per cent of catalogue cores associate with the skeleton (Table~\ref{tab:nnpipeline}), and the rejected cores have projected separations $1.5$--$3.0$ times \emph{larger}, so coverage selection biases the adjacent spacing short. Relaxing the association half-width from $0.06$ to $0.30$\,pc raises coverage to $70$--$97$ per cent and lengthens the median adjacent spacing by only $3$--$16$ per cent (Fig.~\ref{fig:scaletests}c); $S_{\rm local}$ plateaus and the bias is bounded at $\lesssim16$ per cent.

The unconditional statistic behaves differently, and this is a genuine limitation. Because adding cores at fixed $L$ necessarily raises the line density, $S_{\rm global}$ falls monotonically with association radius, from $3.6$ to $2.1$ in Orion~B, $3.6$ to $2.0$ in Aquila, $6.6$ to $4.2$ in Perseus and $5.1$ to $4.0$ in Taurus. The near-classical value is robust in Perseus and Taurus but in Orion~B and Aquila holds only near the adopted radius. The local/global contrast persists in all four clouds; its size in the two low-coverage clouds depends on how generously cores are assigned.

\begin{table}
\caption{Principal statistics under both filament definitions, at a $0.06$\,pc association half-width and a $0.06$\,pc occupancy half-width. $F(1\,{\rm pc})$ is omitted where the FilFinder components are mostly shorter than the cell. The local and global statistics move in opposite senses between the two algorithms, so their contrast is not an artefact of either ontology.}
\label{tab:ontology}
\footnotesize
\setlength{\tabcolsep}{3pt}
\begin{tabular}{llccccc}
\toprule
Region & Skeleton & $L$ & $N_{\rm assoc}$ & $(L/N)/W$ & $s_{\rm ACS,med}/W$ & $f_{\rm occ}$ \\
 & & (pc) & & & & \\
\midrule
Orion B & DisPerSE  & 337 & 995 & 3.4 & 1.77 & 0.27 \\
        & FilFinder & 121 &  73 & 16.6 & 1.54 & 0.06 \\
Aquila  & DisPerSE  & 106 & 308 & 3.4 & 1.66 & 0.27 \\
        & FilFinder &  67 &  40 & 16.7 & 0.97 & 0.05 \\
Perseus & DisPerSE  & 290 & 501 & 5.8 & 2.97 & 0.12 \\
        & FilFinder &  78 &  85 &  9.1 & 1.05 & 0.09 \\
Taurus  & DisPerSE  & 209 & 463 & 4.5 & 1.81 & 0.06 \\
        & FilFinder & 111 & 120 &  9.3 & 0.54 & 0.05 \\
\bottomrule
\end{tabular}
\end{table}

\begin{figure*}
\centering
\includegraphics[width=0.94\textwidth]{figures/fig_scaletests.pdf}
\caption{Scale and convention tests of the inhomogeneity result. \emph{(a)} Occupancy fraction against the adopted half-width $\ell$ (dotted line: the $0.06$\,pc value quoted in the text). Taurus and Perseus saturate, so their low occupancy is not a consequence of the chosen $\ell$; Orion~B and Aquila rise steadily and are less strongly segregated. \emph{(b)} Fano factor against cell length. Over-dispersion ($F>1$, dashed) persists across $0.25$--$4$\,pc in Taurus, Perseus and Orion~B, but Aquila is consistent with Poisson below $1$\,pc. \emph{(c)} Convergence against association half-width: the conditional $s_{\rm ACS,med}/W$ (solid) plateaus, whereas the unconditional $(L/N)/W$ (dashed) falls as more cores are admitted, approaching the classical band (shaded) only in Perseus and Taurus.}
\label{fig:scaletests}
\end{figure*}

\subsection{Spatial inhomogeneity: filament occupancy and scale-dependent over-dispersion}
\label{sec:intermittency}

The gap between $S_{\rm local}$ and $S_{\rm global}$ implies that cores occupy preferential stretches. This can be measured. We assign each associated core an occupancy interval of $\pm0.06$\,pc along its spine, merge overlapping intervals, and define the occupancy fraction $f_{\rm occ}$ as the fraction of skeleton length that is core-bearing; we also count cores in $1$\,pc cells and form the Fano factor $F=\sigma^2/\mu$ and the Gini coefficient of those counts (Table~\ref{tab:intermittency}, Fig.~\ref{fig:intermittency}).

On this operational definition between $6$ and $27$ per cent of the skeleton is core-bearing, so most filament length is empty at HGBS sensitivity. That number is tied to the adopted $\pm0.06$\,pc interval and is not an independently measured filling factor; the scale-independent evidence for inhomogeneity is the Fano behaviour and the length distribution of core-free runs, given below. The Fano factor exceeds unity everywhere ($1.6$, $2.8$, $5.7$, $25.2$ for Aquila, Orion~B, Perseus and Taurus) and the Gini coefficients are $0.36$--$0.57$. In Orion~B, $92$\,pc of the $337$\,pc skeleton is occupied, a further $155$\,pc lies in core-free runs inside core-bearing components, and roughly $90$\,pc is in components with no catalogued core at all.

Because $f_{\rm occ}$ necessarily grows with the assumed interval, we report it against $\ell$ (Fig.~\ref{fig:scaletests}a). Taurus rises only from $0.052$ to $0.078$ and Perseus from $0.078$ to $0.202$ as $\ell$ increases tenfold, so widening the interval adds almost nothing; Orion~B and Aquila climb to $\approx0.6$ and are much less segregated. The Fano factor is the more robust statistic because it needs no occupancy convention: $F(\ell)>1$ across $\ell=0.25$--$4$\,pc in Taurus ($9.4$--$25$), Perseus ($2.3$--$8.4$) and Orion~B ($1.2$--$4.9$), rising monotonically with $\ell$. Aquila is the exception, consistent with Poisson below $1$\,pc.

\smallskip\noindent\textbf{Is this more than the filament's own structure implies?} A Fano factor above unity compares the cores with a \emph{homogeneous} Poisson process, and a filament is not homogeneous. We therefore built an inhomogeneous null whose intensity is conditioned on the local column density, measured $\lambda(N_{\rm H_2})$ from the data, scattered synthetic cores along the actual skeletons and passed them through the same pipeline, in $10^4$ realisations per cloud (Appendix~\ref{sec:nullspec}). In Orion~B, Aquila and Perseus the observed $F$ lies \emph{below} the null median ($p=0.98$, $0.88$, $0.96$): the column-density structure alone accounts for all of the over-dispersion. Only Taurus exceeds its null ($F=25.2$ against a median of $13.4$; $p=2\times10^{-4}$). Refitting the intensity out of sample, from the other three clouds, leaves the conclusion unchanged: $p_{\rm cv}=1.00$, $0.27$, $0.995$ and $0.005$ for Orion~B, Aquila, Perseus and Taurus. Both versions are given in Table~\ref{tab:percloud} and we quote $p_{\rm fit}$ in the text unless stated otherwise.

One bias remains and it works against us: cores raise the column density of their own surroundings, so the pixels to which the null assigns the highest intensity are the core sites themselves, and the null is too easily satisfied. Masking a $0.06$\,pc radius around each associated core changes the intensity in the highest column bins by factors of $4$--$10$. The three non-detections are therefore conservative rather than secure, and the Taurus excess is correspondingly strengthened. Excess clumping is therefore detected in one of the four clouds and not detected in the other three, which is not the same as its being absent there.

Two length scales coexist. On sub-parsec scales the gap distribution shows a deficit of the smallest separations (Section~\ref{sec:clustering}); on parsec scales the counts are strongly over-dispersed. Cores form in groups with short internal separations, and the groups are separated by long stretches that have produced none. This accounts for the factor of about two between $S_{\rm local}$ and $S_{\rm global}$ without appealing to small-scale clustering, and it is the central observational result of this paper: what a fragmentation model must reproduce is a spatial point process, not a single spacing.

\begin{table}
\caption{Spatial-inhomogeneity diagnostics. $L_{\rm tot}$ is the total skeleton length, $L_{\rm core}$ the core-bearing length (union of $\pm0.06$\,pc intervals about each associated core), and $f_{\rm occ}$ is their ratio. $f_{\rm assoc}$ is the fraction of catalogue cores associated with the skeleton, $F$ the Fano factor of counts in $1$\,pc cells and $G$ their Gini coefficient. Both association half-widths are listed. }
\label{tab:intermittency}
\footnotesize
\setlength{\tabcolsep}{3pt}
\begin{tabular}{lccccccccc}
\toprule
Region & half & $L_{\rm tot}$ & $L_{\rm core}$ & $f_{\rm occ}$ & $f_{\rm assoc}$ & $\frac{L/N}{W}$ & $\frac{L_{\rm core}/N}{W}$ & $F$ & $G$ \\
       & (pc) & (pc) & (pc) & & & & & & \\
\midrule
Orion B & 0.06 & 337 & 92 & 0.27 & 0.53 & 3.4 & 0.9 & 2.8 & 0.49 \\
        & 0.30 & 337 & 117 & 0.35 & 0.85 & 2.1 & 0.7 & 4.5 & 0.50 \\
Aquila  & 0.06 & 106 & 29 & 0.27 & 0.41 & 3.4 & 0.9 & 1.6 & 0.40 \\
        & 0.30 & 106 & 38 & 0.36 & 0.70 & 2.0 & 0.7 & 2.1 & 0.36 \\
Perseus & 0.06 & 290 & 35 & 0.12 & 0.61 & 5.8 & 0.7 & 5.7 & 0.43 \\
        & 0.30 & 290 & 40 & 0.14 & 0.85 & 4.2 & 0.6 & 10.5 & 0.47 \\
Taurus  & 0.06 & 210 & 13 & 0.06 & 0.86 & 4.5 & 0.3 & 25.2 & 0.57 \\
        & 0.30 & 210 & 14 & 0.07 & 0.97 & 4.0 & 0.3 & 27.3 & 0.57 \\
\bottomrule
\end{tabular}
\end{table}

\begin{figure}
\centering
\includegraphics[width=0.99\columnwidth]{figures/fig_intermittency.pdf}
\caption{Spatial inhomogeneity of core occurrence along the filaments. \emph{(a)} The $0.06$\,pc occupancy fraction (bars, left axis) is $6$--$27$\,per cent, while the Fano factor of counts in $1$\,pc cells (red, right axis, logarithmic) exceeds the Poisson value of unity in every cloud. \emph{(b)} The three normalised spacing statistics: the local adjacent-core spacing is well below the classical band, the global $(L/N)/W$ approaches it, and the core-bearing $(L_{\rm core}/N)/W$ lies below both because the occupied stretches are densely populated. The separation between the blue and green curves measures how unevenly the cores occupy the crest.}
\label{fig:intermittency}
\end{figure}

\subsection{Regularity of the core spacing}
\label{sec:clustering}

Classical fragmentation predicts a quasi-periodic chain; fibre projection or a turbulent cascade predicts a broad distribution with no comb. Ordering associated cores by arclength and measuring adjacent gaps gives a coefficient of variation of order unity (Taurus $1.36$, Orion~B $0.83$, Perseus $0.96$, Aquila $0.65$; Fig.~\ref{fig:clustering}).

Goodness-of-fit tests were calibrated by parametric bootstrap ($N_{\rm boot}=10^5$), resampling from the fitted distribution to build the exact null distribution of the Kolmogorov--Smirnov statistic. The result is two-sided: the adopted periodic-comb-plus-jitter model is rejected ($p_{\rm boot}<10^{-5}$ in every cloud, this being the floor set by the number of resamples rather than a measured value) \emph{and} so is the pure exponential ($p_{\rm boot}<10^{-5}$ in three clouds, $1.7\times10^{-3}$ in Perseus). Comparing five renewal models by the corrected Akaike criterion gives $\Delta$AICc$=37$--$246$ against Poisson and $81$--$742$ against the periodic comb in every cloud. The shifted-exponential hard-core model is not uniquely preferred once broader alternatives are admitted, a positively-skewed lognormal fitting best in three clouds, so the data favour a broad, skewed distribution truncated at small separations rather than a sharp physical hard core.

These are, strictly, rejections of two specific nulls. Because a real fragmenting filament would not produce an ideal comb, we therefore tested a deliberately imperfect one: a chain whose local wavelength varies slowly along the crest, with positional jitter and a random fraction of fragments removed. Two versions were tried, a mild case ($\pm25$ per cent wavelength variation, $25$ per cent jitter, $15$ per cent dropout) and a strong case ($\pm50$ per cent, $25$ per cent, $35$ per cent). Both are rejected in all four clouds at the bootstrap floor ($p_{\rm boot}=3\times10^{-3}$, set by $300$ resamples), with Kolmogorov--Smirnov distances of $0.13$--$0.23$. The observed gap distributions are therefore inconsistent not only with an idealised comb but with a physically degraded one, which is a stronger statement than the AICc comparison alone supports. What remains unexcluded is a chain whose wavelength varies systematically with local filament properties rather than randomly, which would need those properties measured per stretch.

The clouds fail the exponential for different reasons and should not be summarised together (Table~\ref{tab:percloud}). Writing $R$ for the observed median gap divided by that expected for an exponential of the same mean, only Orion~B ($1.03$) and Aquila ($1.14$) show a deficit of the smallest separations; Perseus ($0.92$) and Taurus ($0.64$) show an excess.

\smallskip\noindent\textbf{The small-separation deficit.} The gap distributions are truncated below $0.011$--$0.044$\,pc. This range lies entirely inside the extraction floor: the $18.2''$ beam corresponds to $0.012$--$0.039$\,pc across our distance range and \textsc{getsources} does not reliably separate sources closer than about one beam. We therefore treat the truncation as observational and attach no physical interpretation to it. Deciding whether any physical suppression is also present requires a completeness function $C(s,F_1/F_2,N_{\rm bg})$ measured by injecting synthetic source pairs of varying separation, flux ratio, size and local background into the maps and re-extracting them through the HGBS pipeline, so that the gap distribution can be forward-modelled rather than truncated. That experiment is not attempted here.

\smallskip\noindent\textbf{Gap independence.} A runs test about the median gap shows significant non-independence in Perseus ($z=-4.2$) and Taurus ($z=-5.4$), the same core-poor/core-rich alternation documented in Section~\ref{sec:intermittency}. We therefore recomputed $\Delta$AICc with a moving-block bootstrap resampling contiguous runs of five adjacent gaps drawn from within single components ($400$ resamples per cloud). The effective sample size falls to $122$ of $440$ gaps in Taurus and $344$ of $438$ in Perseus, but the rejections survive: the $5$th percentile of the bootstrap $\Delta$AICc is $415$, $45$, $249$ and $587$ against the periodic model and $190$, $71$, $31$ and $55$ against Poisson. What does not survive is discrimination among the broad models, the shifted-exponential lying inside the lognormal's bootstrap interval in three clouds.

\smallskip\noindent\textbf{No classical-scale periodicity is detected.} The along-skeleton pair-correlation $g(s)$ (Fig.~\ref{fig:paircorr}) shows no feature at $5$--$6\,W$. To decide whether that absence is significant we generated a $95$ per cent envelope from the column-density-conditioned null of Appendix~\ref{sec:nullspec}: at the classical scale the observed $g$ lies \emph{inside} the envelope in every cloud (Orion~B $1.02$ in $[0.84,1.26]$; Aquila $0.89$ in $[0.58,1.52]$; Perseus $1.02$ in $[0.91,1.14]$; Taurus $1.00$ in $[0.91,1.10]$). This is a non-detection at the $95$ per cent level rather than a demonstration of absence, and it is more informative than a comparison of median spacings because it does not compress the point process into a single number.

\begin{figure}
\centering
\includegraphics[width=0.99\columnwidth]{figures/fig_paircorr.pdf}
\caption{Along-skeleton pair-correlation function $g(s)$ for the four primary clouds (uniform-random expectation on the matched components divided out; $g=1$ is no correlation). Dotted vertical lines mark the classical scale $5.5\,W_{\rm fil}$ per region. No cloud shows a peak at the classical scale ($g\approx1$ there, slightly below unity in Orion~B and Aquila); the only departures are a mild small-separation excess (Taurus) and weak region-envelope features near $1$\,pc. The classical periodic comb is therefore absent, consistent with a non-periodic, spatially inhomogeneous distribution bearing an observational small-separation deficit (Section~\ref{sec:clustering}).}
\label{fig:paircorr}
\end{figure}

\begin{figure}
\centering
\includegraphics[width=0.99\columnwidth]{figures/fig_clustering.pdf}
\caption{Adjacent-core gap distributions along the skeletons of the four primary regions (histograms), with exponential/Poisson (solid) and regular/periodic (dashed) reference curves; The shaded region at small separations spans the \textit{Herschel} beam and \textsc{getsources} deblending scale for each cloud; \emph{the distribution below approximately one beam is not interpreted astrophysically}. Each panel gives the coefficient of variation CoV, the 1-D Clark--Evans ratio $R$, and the exponential goodness-of-fit $p$-value \emph{calibrated by parametric bootstrap} ($p_{\rm boot}$; the null distribution of the KS statistic is built by resampling from the fitted exponential, removing the estimated-parameter bias of the standard test). The result is two-sided: the regular/periodic model is rejected ($p_{\rm boot}<10^{-5}$, the parametric-bootstrap floor, $1/N_{\rm boot}$ with $N_{\rm boot}=10^5$) in every region, \emph{and} the pure-exponential (Poisson) model is also rejected ($p_{\rm boot}<10^{-5}$ (Perseus $1.7\times10^{-3}$) in all four), so in each of the four clouds the cores are neither periodic nor pure-Poisson; consistent with a point process that is limited at small separations by the extraction resolution but over-dispersed on parsec scales (the deficit of small gaps, $R>1$; largely instrumental).}
\label{fig:clustering}
\end{figure}

\begin{table}
\caption{Cloud-by-cloud summary of the point-process diagnostics, to avoid summarising four different behaviours as one. $R$ is the observed median gap divided by the median expected for an exponential of the same mean ($R>1$ deficit of small gaps, $R<1$ excess); $F_{\rm obs}$ is the observed Fano factor of counts in $1$\,pc cells and median $F_{\rm null}$ that of the inhomogeneous, column-density-conditioned null of Appendix~\ref{sec:nullspec}. Two $p$-values are given because two versions of the null exist and both are reported throughout: $p_{\rm fit}$ uses an intensity exponent fitted to the cloud's own cores, and $p_{\rm cv}$ one predicted from the other three clouds under leave-one-cloud-out cross-validation. Both are empirical one-sided values from $10^4$ realisations. The two agree except in Aquila, whose fitted exponent is the least well constrained; Taurus exceeds its null on either version. The clouds do not behave alike, and only Taurus exceeds its inhomogeneous null. With $N=4$ this is a single case, not a demonstrated population split; and because the null is built partly from column density the cores themselves raise, the three non-detections are conservative.}
\label{tab:percloud}
\footnotesize
\setlength{\tabcolsep}{3.5pt}
\begin{tabular}{lcccccc}
\toprule
Region & CoV & $R$ & $F_{\rm obs}$ & median $F_{\rm null}$ & $p_{\rm fit}$ & $p_{\rm cv}$ \\
\midrule
Orion B & 0.83 & 1.03 & 2.8  & 3.6  & 0.98 & 1.00 \\
Aquila  & 0.65 & 1.14 & 1.6  & 2.1  & 0.88 & 0.27 \\
Perseus & 0.96 & 0.92 & 5.7  & 7.1  & 0.96 & 0.995 \\
Taurus  & 1.36 & 0.64 & 25.2 & 13.4 & $\mathbf{2\times10^{-4}}$ & $\mathbf{0.005}$ \\
\bottomrule
\end{tabular}
\end{table}

\subsection{Systematic error budget}
\label{sec:error_budget}

Not all of the checks reported above carry equal weight, and it is worth saying which are decisive. With $N=4$ clouds, the tests that bear directly on the conclusions are those that could change the \emph{sign} of a result: the skeleton-ontology comparison (Section~\ref{sec:ontology}), the coverage and association-radius sweeps, the core-population splits (Section~\ref{sec:popsplit}) and the column-density-conditioned null (Section~\ref{sec:intermittency}). The remainder, including the width and persistence conventions, the injection--recovery corrections and the eligible-length sweeps, are illustrative: they move numbers within the envelope already quoted and none of them reverses a conclusion. Table~\ref{tab:errorbudget} collects the terms. We separate \emph{statistical/sensitivity} terms, which would change under resampling, from \emph{systematic} offsets and from \emph{modelling} terms arising in the simulations. Uncorrelated statistical terms are combined in quadrature, $\sigma_{\rm sys}^2=\sum_i\sigma_i^2$; the multiplicative width- and skeleton-convention factors are applied afterwards as separate analysis-choice envelopes rather than added into that sum. The adopted $(s_{\rm ACS,med}/W)_0\approx1.4$ carries residual pipeline bias ($\sim10$ per cent), protostellar migration ($\sim7.5$ per cent), between-region scatter ($\sim4$ per cent) and a $\lesssim5$ per cent manual-editing proxy in quadrature, with the width convention ($\sim25$ per cent) and the skeleton algorithm (a factor $\sim2$) applied on top. Because the dominant term is a deterministic offset we quote a min--max envelope, $S_{\rm local}\approx0.8$--$1.9$ for DisPerSE, rather than a Gaussian $\sigma$.

\emph{Projection} acts against our conclusion and is shown in Table~\ref{tab:headline} for that reason. For an isotropic orientation distribution $s_{\rm obs}/s_{\rm true}=\sin i$ has median $0.87$, so deprojection raises the intrinsic spacing by $15$--$27$ per cent, moving the adopted $1.4$ to $1.6$--$1.8$. We do not apply it, because the orientation distribution within a single cloud need not be isotropic, but we carry it as a one-sided systematic.

\emph{Skeleton-algorithm dependence} is the largest observational term. The DisPerSE and FilFinder values ($\approx1.9$ and $\approx0.4$--$0.8$) bracket a genuine ambiguity in what constitutes ``the filament'' rather than a measurement error, and a matched-threshold FilFinder extraction shows the inferred spacing is itself strongly sensitive to the mask threshold, so the factor of two is if anything a lower bound. We base our quantitative conclusions on the persistent-crest network, which is the physically better-motivated definition of a filament ridge, and treat FilFinder as an ontology test rather than an equally weighted measurement.

\begin{table}
\caption{Summary of the terms on the $\lambda/W$ comparison (symbols: Table~\ref{tab:symbols}), grouped as \emph{observational} (on $s_{\rm ACS,med}$), \emph{simulation} (on $\lambda_{\rm turn}$) and \emph{comparison} (terms that enter only when the two are matched). Statistical/sensitivity terms are combined in quadrature; the width- and skeleton-convention factors are applied afterwards as a separate multiplicative \emph{analysis-choice envelope}, not as Gaussian uncertainties.}
\label{tab:errorbudget}
\footnotesize
\setlength{\tabcolsep}{3pt}
\begin{tabular}{@{}p{3.1cm}p{1.9cm}p{2.3cm}@{}}
\toprule
Source & Type & Magnitude \\
\midrule
\multicolumn{3}{@{}l}{\emph{Observational (on $s_{\rm ACS,med}$)}}\\
Bootstrap / cloud-to-cloud & sensitivity & $\sigma\approx0.4$ ($\sim$20\%) \\
ACS pipeline bias & systematic & $\sim$10\% (corrected) \\
Protostellar migration & sensitivity & $\sim$7.5\% \\
Projection (deprojection) & systematic (one-sided, \emph{raises} $s$) & $+15$--$27\%$ \\
LOS blending (Aquila) & systematic (one-sided, \emph{shortens} $s$) & $\lesssim50\%$ worst case \\
Skeleton-coverage selection & systematic (one-sided, \emph{shortens} $s$) & $\lesssim16\%$ (bounded) \\
Width convention, fixed vs region & analysis-choice env. & $\sim$25\% \\
Skeleton algorithm, DisPerSE/FilFinder & analysis-choice env. & factor $\sim$2 \\
Manual skeleton editing & systematic (proxy) & $\lesssim$5\% \\
\midrule
\multicolumn{3}{@{}l}{\emph{Simulation (on $\lambda_{\rm turn}$)}}\\
Growth-rate resolution/epoch & modelling & $\lesssim$5\% (converged) \\
Turbulence sensitivity & sensitivity & $<$5\% (single realisation) \\
Line-mass ($f$) dependence & modelling & $\sim$25\% \\
\midrule
\multicolumn{3}{@{}l}{\emph{Comparison (matching the two)}}\\
Gaussian vs Plummer FWHM & systematic & factor $\sim$1.5 \\
$T_1$ width normalisation ($\eta_W$) & systematic & $\sim$30\% \\
\bottomrule
\end{tabular}
\end{table}

\subsection{Pipeline counts, coverage and population splits}
\label{sec:nnpipeline}

The adjacent-core measurement thresholds the deposited DisPerSE skeleton, bridges fragmented segments by morphological closing and re-skeletonisation, associates each core to the nearest spine within $0.06$\,pc, orders the associated cores along each component by graph arclength (double-BFS diameter) and takes adjacent separations. Table~\ref{tab:nnpipeline} lists the resulting counts.

\begin{table}
\caption{Measurement counts per primary region: catalogue cores, cores associated to a skeleton within $0.06$\,pc, the associated fraction, and the number of adjacent-core pairs entering the median.}
\label{tab:nnpipeline}
\footnotesize
\setlength{\tabcolsep}{5pt}
\begin{tabular}{lcccc}
\toprule
Region & catalogue & associated & $f_{\rm assoc}$ & ACS pairs \\
\midrule
Orion B & 1870 & 995 & 0.53 & 699 \\
Aquila  &  749 & 308 & 0.41 & 168 \\
Perseus &  816 & 501 & 0.61 & 438 \\
Taurus  &  536 & 463 & 0.86 & 440 \\
\bottomrule
\end{tabular}
\end{table}

The associated fraction ranges from $0.41$ to $0.86$, so in Orion~B and Aquila the spacing is measured on roughly half the catalogue. This selection is not neutral, but it acts \emph{towards} the sub-classical result: the unassociated cores lie on fainter segments and have projected separations $1.5$--$3.0$ times larger, and relaxing the association half-width to $0.30$\,pc raises the coverage to $70$--$97$ per cent while lengthening the median adjacent spacing by only $3$--$16$ per cent (Section~\ref{sec:ontology}). The bias on $S_{\rm local}$ is therefore one-sided and bounded at $\lesssim16$ per cent. Sweeping the half-width over $0.04$--$0.08$\,pc changes the per-region median by $4$--$7$ per cent, so the result is not sensitive to that choice.

Restricting to the two clouds with the highest coverage, Perseus and Taurus, gives the same median as the full four-region sample. The dependence on the core population is examined in Section~\ref{sec:popsplit}.


\subsection{Core population splits}
\label{sec:popsplit}

The HGBS catalogues contain unbound starless cores, prestellar cores and protostellar sources. These are not equivalent for the present purpose: protostellar objects have had time to accrete, migrate and interact, while unbound starless cores need not be fragmentation products at all. We therefore repeated the principal statistics on three nested samples per cloud (Table~\ref{tab:popsplit}).

The conditional statistic is essentially invariant. $S_{\rm local}$ changes by at most $0.2$ between the all-core and prestellar-only samples in every cloud, and the coefficient of variation of the gaps is unchanged to within $0.1$. The unconditional statistic behaves in the opposite way, rising steeply as the sample is restricted, from $3.4$ to $5.9$ in Orion~B, $3.4$ to $4.8$ in Aquila and $5.8$ to $10.3$ in Perseus, simply because $L$ is fixed while $N$ falls.

That combination sharpens rather than weakens the central result. The separation between the local and global statistics is not produced by the inclusion of evolved or unbound objects; restricting to the cores most likely to be genuine fragmentation products \emph{widens} it, and pushes $S_{\rm global}$ from near-classical to super-classical in three of the four clouds. It also shows that $S_{\rm global}$ is sensitive to catalogue selection in a way that $S_{\rm local}$ is not, which is a further reason to treat the comparison of $S_{\rm global}$ with $\lambda_{\rm m}$ as a diagnostic rather than a measurement. The Taurus prestellar sample contains only $54$ cores and its $S_{\rm global}$ is not meaningful; we quote it only for completeness.

\begin{table}
\caption{Principal statistics for three nested core samples. ``All'' is the full published catalogue; ``starless$+$prestellar'' removes protostellar sources; ``prestellar'' retains only cores classified as prestellar. $N_{\rm assoc}$ is the number associated with a skeleton, CoV the coefficient of variation of the adjacent gaps. The Taurus prestellar row rests on $53$ associated cores and its $S_{\rm global}$ should not be interpreted.}
\label{tab:popsplit}
\footnotesize
\setlength{\tabcolsep}{4pt}
\begin{tabular}{llcccc}
\toprule
Cloud & sample & $N_{\rm assoc}$ & $S_{\rm local}$ & $S_{\rm global}$ & CoV \\
\midrule
Orion~B & all                  & 995 & 1.77 & 3.4 & 0.83 \\
        & starless$+$prestellar& 961 & 1.79 & 3.5 & 0.83 \\
        & prestellar           & 569 & 1.83 & 5.9 & 0.82 \\
\midrule
Aquila  & all                  & 308 & 1.66 & 3.4 & 0.65 \\
        & starless$+$prestellar& 267 & 1.78 & 4.0 & 0.62 \\
        & prestellar           & 223 & 1.78 & 4.8 & 0.73 \\
\midrule
Perseus & all                  & 501 & 2.97 & 5.8 & 0.96 \\
        & starless$+$prestellar& 438 & 2.99 & 6.6 & 0.95 \\
        & prestellar           & 282 & 2.83 & 10.3 & 1.01 \\
\midrule
Taurus  & all                  & 463 & 1.81 & 4.5 & 1.36 \\
        & starless$+$prestellar& 424 & 1.84 & 4.9 & 1.34 \\
        & prestellar           &  53 & 2.04 & (39.5) & 1.30 \\
\bottomrule
\end{tabular}
\end{table}

\subsection{Comparison with the literature}

Our revised-distance ACS measurements (Table~\ref{tab:literature}) differ from the original HGBS values by the region-dependent distance revisions; the core-count-weighted primary-region mean rises only $\sim10$ per cent, dominated by Aquila. Toggling each methodological choice on the Orion~B data at fixed distance shows that the \emph{statistic definition} dominates: switching from ACS to the pairwise median inflates the spacing from $0.207$ to $0.313$\,pc and accounts for essentially the whole apparent shift, while skeleton re-extraction, association half-width, arclength ordering and the distance revision are each $\leq15$ per cent and largely cancel. On our primary statistic Orion~B agrees with \citet{Konyves2020} to within $3$ per cent.

The published values we could check are local or along-fibre separations, not all-pairs medians, which is why they agree with our ACS column rather than our pairwise column. The region-extent bias applies wherever an all-pairs or pooled statistic is adopted, including the comparability statistic carried here and in cross-region compilations, and we make no claim that it affects these particular published values.

\begin{table*}
\caption{Comparison with published HGBS spacing measurements. Perseus and Taurus are the two regions in which our own gap sequence shows significant non-independence (runs test; Section~\ref{sec:clustering}) and in which the post-masking eligible-core samples collapse (Section~\ref{sec:eligible}); their entries should be compared with the published values with those caveats in mind.}
\label{tab:literature}
\setlength{\tabcolsep}{4pt}
\footnotesize
\setlength{\tabcolsep}{2.5pt}
\begin{tabular}{lcccccc}
\toprule
Region & ACS (pc) & Orig.\ HGBS & pairwise (pc) & Lit.\ (pc) & \textbf{Lit.\ statistic type} & Reference \\
\midrule
\textbf{Primary sample} &  &  &  &  &  &  \\
Orion B & \textbf{0.207} & 0.212 & 0.313 & 0.20--0.30 & local & \citet{Konyves2020} \\
Aquila  & \textbf{0.210} & 0.206 & 0.346 & 0.22--0.26 & local & \citet{Konyves2015} \\
Perseus & \textbf{0.263} & 0.199 & 0.248 & 0.22 & local & \citet{Pezzuto2021} \\
Taurus  & \textbf{0.124} & 0.206 & 0.198 & 0.19 & along-fibre & \citet{Hacar2013} \\
\midrule
\textbf{Limited sample} &  &  &  &  &  &  \\
Ophiuchus &, & 0.197 & 0.206 & 0.18 & local & \citet{Konyves2020} \\
Serpens  &, & 0.188 & 0.331 & 0.17--0.19 & local & \citet{Konyves2020} \\
TMC1     &, & 0.203 & 0.195 & $\sim$0.20 & along-fibre & \citet{Hacar2013} \\
CRA      &, & 0.212 & 0.248 & $\sim$0.21 & local & \citet{Konyves2020} \\
\bottomrule
\end{tabular}

\vspace{0.1in}
\footnotesize
\emph{Notes} $^{a}$ The ``Lit.\ statistic'' column records what kind of statistic each cited value represents. All are \emph{local} separation measures, either nearest-neighbour or along-fibre adjacent separations, rather than all-pairs medians, and that is precisely why they agree with our ACS column rather than with our pairwise column; the exact definitions differ in detail between papers. The region-extent bias identified here therefore applies wherever an all-pairs or pooled separation statistic is adopted, including the comparability statistic carried in this paper and in cross-region compilations, and not to these particular published values. $^{b}$ The primary statistic is the \emph{adjacent-core} (ACS) spacing (bold); the ``pairwise'' column is a legacy comparator only, inflated towards $L/3$ (Section~\ref{sec:statistics}) and not a physical spacing. ACS spacings for the limited sample are omitted. Perseus is the one region where the ACS ($0.263$) slightly exceeds the pairwise median, because arclength ordering on its heavily bridged skeleton inflates path distances; we verified this is not a bridging artefact, an independent minimum-spanning-tree ordering of the same cores gives a consistent value, insensitive to the morphological bridging, and both remain sub-classical. $^{c}$ ``Original HGBS'' values use pre-Gaia distances; literature ranges reflect different methods and distance assumptions. $^{d}$ For Orion~B the pre-Gaia distance was already $\sim$400\,pc \citep{Konyves2020}, so the pairwise $0.212\to0.313$\,pc difference is a statistic-definition effect, not a distance revision.
\end{table*}

\subsection{What constitutes the theoretical filament?}
\label{sec:whatfilament}

The factor-of-two DisPerSE-versus-FilFinder difference (Fig.~\ref{fig:skeleton_compare}) is not measurement noise: the two algorithms \emph{define different objects} (persistent crests versus the medial axis of a masked region), so there is no single, algorithm-independent ``filament'' whose spacing they should reproduce. Our robustness claim applies to the \emph{sign} of the deficit, every convention giving a sub-classical $S_{\rm local}$, while its \emph{magnitude} cannot be assigned a unique value: the same data yield $S_{\rm local}\approx1.0$ (region-Plummer FilFinder), $\approx1.4$ (region-Plummer DisPerSE) or $\approx1.9$ (fixed-width DisPerSE). Because segmentation alone moves the answer by as much as the differences between the physical models compared in Section~\ref{sec:validation_endtoend}, it should not be read as discriminating between them. Figure~\ref{fig:envelope} makes the resulting envelope explicit.

\begin{figure}
\centering
\includegraphics[width=0.99\columnwidth]{figures/fig_envelope.pdf}
\caption{The full methodological envelope of $\lambda/W$, organised by \emph{which quantity is being measured} rather than by convention alone. The four groups are the conditional adjacent-core spacing $S_{\rm local}$ (blue), the unconditional length per core $S_{\rm global}=(L/N)/W$ (green), the same restricted to fragmentation-eligible filament $S_{\rm elig}$ (purple; the core-masked row is fiducial, the raw-map row shows the endogeneity-affected comparison), and the simulated upper bound $S_{\rm sim}^{\rm upper}=\lambda_{\rm turn}/W$ together with the hierarchical-fibre toy model (red). Each bar spans the range across clouds and conventions; the grey band is the classical $\lambda_{\rm m}\approx5$--$6$ (\citealt{Inutsuka1992,Fischera2012}) and the dotted line the critical $\lambda_{\rm cr}\approx2.6$. Grouping this way makes the central point of the paper visible directly: $S_{\rm local}$ lies well below the classical band under every convention, $S_{\rm global}$ reaches it, $S_{\rm elig}$ falls between the two, and $S_{\rm sim}^{\rm upper}$ is an \emph{upper bound} that should not be read as a measured wavelength. Comparisons between groups are comparisons between different quantities and are only meaningful once that is stated (Section~\ref{sec:whatfilament}). The horizontal axis is symmetric-log below $\lambda/W=1$. The $S_{\rm sim}^{\rm upper}$ bar is drawn with a left-pointing arrow because it is an upper limit, not a measured interval; the open marker on the $S_{\rm local}$ group shows where the adopted value moves if the isotropic deprojection factor of Table~\ref{tab:headline} is applied.}
\label{fig:envelope}
\end{figure}

There is a more fundamental point behind this. The classical calculation is an eigenmode of a single, coherent, approximately cylindrical, self-gravitating object. Neither a persistent column-density crest nor the medial axis of a mask guarantees that such an object exists along the traced path: both are two-dimensional constructions on a projected scalar field, and either may connect material that is unrelated in velocity, or divide material that is dynamically single. Before any skeleton-derived length is compared with a cylinder eigenmode, the segment must be shown to correspond to a velocity-coherent structure. That cannot be established from column density, and it is the reason the comparison in this paper is presented as a diagnostic rather than a test, and the reason the decisive follow-up is kinematic: dense-gas position--position--velocity data that identify which of the two networks, if either, traces a dynamically coherent cylinder.

\section{Theoretical framework}
\label{sec:theory}

\subsection{Classical fragmentation theory}




For an isothermal, self-gravitating cylinder of density $\rho_0$ and sound speed $c_s$, the critical line mass for gravitational instability is \citep{Ostriker1964}:
\begin{equation}
\mu_{\rm crit} = \frac{2c_s^2}{G}
\label{eq:mucrit}
\end{equation}
The most unstable fragmentation wavelength is, as an \emph{exact} result for the isothermal equilibrium cylinder \citep{Inutsuka1992},
\begin{equation}
\lambda_{\rm IM92} = 22\,H \approx 4\times(2R_{\rm flat}),
\label{eq:lambdam}
\end{equation}
where $H = c_s\,(4\pi G\rho_c)^{-1/2}$ is the thermal scale height and $R_{\rm flat}=\sqrt8\,H$ the flattening radius (\citealt{Nagasawa1987}, eq.~33; \citealt{Fischera2012}, eqs.~5 and 24; verified independently to better than $1$ per cent in Appendix~\ref{sec:validation}), and $W_{\rm fil}\approx0.1$\,pc the characteristic filament FWHM. Solving the dispersion relation exactly (Appendix~\ref{sec:validation}; \citealt{Fischera2012}) gives a fastest-growing wavelength $\lambda_{\rm m}\approx5\times W_{\rm fil}$ near-critical, rising to $\approx6\times W_{\rm fil}$ at low line mass, and a critical wavelength $\lambda_{\rm cr}\approx2.6\times W_{\rm fil}$. The often-quoted ``$\lambda\approx4\times$ the width'' is $4\times$ the flat \emph{diameter} $2R_{\rm flat}$ \citep{Inutsuka1992}, which equals $\approx5$--$6\times$ the \emph{FWHM}; taking it as $4\times$ the FWHM understates the classical scale. We therefore adopt $\lambda_{\rm m}\approx5$--$6\,W_{\rm fil}$ (Fischera \& Martin 2012) as the benchmark. Because we compare in FWHM units and forward-model the simulations to a beam-convolved Plummer FWHM (Section~\ref{sec:width_correction}), the sub-classical deficit does not depend on the equilibrium-vs-observed profile-index difference ($p=4$ vs $p\simeq2$), and we compare as the density-independent ratio $\lambda/W_{\rm fil}$. The deficit is not new in sign: \citet{Fischera2012} noted observed cores sit ``about the FWHM'' apart, and \citet{Zhang2023cal} find a spacing $\approx1.3\times$ the width in California.

\begin{table}
\caption{Where the $5$--$6\times$FWHM benchmark comes from, for the isothermal equilibrium (Ostriker, $p=4$) cylinder. The ratios in the last column are exact for that equilibrium; the final row is the one that is \emph{not} invariant when the observed profile has $p\simeq2$.}
\label{tab:benchmark}
\footnotesize
\setlength{\tabcolsep}{4pt}
\begin{tabular}{llc}
\toprule
Quantity & Definition & Value \\
\midrule
$H$ & $c_s(4\pi G\rho_c)^{-1/2}$, scale height & $1$ \\
$R_{\rm flat}$ & flattening radius, $\sqrt{8}\,H$ & $2.83\,H$ \\
$2R_{\rm flat}$ & flat diameter & $5.66\,H$ \\
$W_{\rm FWHM}$ & FWHM of the $p=4$ profile & $\approx3.9\,H$ \\
$\lambda_{\rm cr}$ & cutoff wavelength & $\approx10\,H\;(\approx2.6\,W)$ \\
$\lambda_{\rm m}$ & fastest-growing wavelength, $22H$ & $22\,H\;(\approx5.6\,W)$ \\
\midrule
$\lambda_{\rm m}/(2R_{\rm flat})$ & the ``$4\times$ diameter'' form & $3.9$ \\
$\lambda_{\rm m}/W_{\rm FWHM}$ & the benchmark used here & $5$--$6$ \\
\bottomrule
\end{tabular}
\end{table}

Two things follow. The ratio $\lambda_{\rm m}/H=22$ and the conversion $\lambda_{\rm m}\approx4\times(2R_{\rm flat})$ are properties of the $p=4$ equilibrium and are exact within it. What is \emph{not} invariant is $\lambda_{\rm m}/W_{\rm FWHM}$, because the FWHM of a $p\simeq2$ Plummer profile is a different fraction of $H$ from that of a $p=4$ profile, and beam convolution broadens it further. The observational benchmark is therefore convention-dependent at the tens-of-per-cent level even before any measurement error, which is one reason we treat the simulated comparison through forward modelling (Section~\ref{sec:width_correction}) rather than through this ratio.

With $f=\mu_{\rm line}/\mu_{\rm crit}$ and $t_{\rm frag}\propto(f-1)^{-1/2}$, marginally supercritical filaments fragment slowly. At magnetically supercritical field strength ($\beta\gtrsim0.2$) all configurations fragment for $\mathcal{M}\geq1$ at every $f$ tested once integrated adequately: above the critical line mass there is no \emph{thermal} stability boundary. This is distinct from \emph{magnetic} stabilisation, where a strong longitudinal field ($\beta\lesssim0.15$) suppresses collapse even for $f>1$ (the ``strong-field, magnetically supported'' regime; the deposited material); the thermal ($f$) and magnetic (mass-to-flux) criticalities are physically distinct.

We adopt a single reference $W_{\rm fil}=0.10$\,pc for comparability with prior HGBS work, though its universality is debated \citep{Panopoulou2017,Panopoulou2022}; the directly-measured region-Plummer widths ($0.11$--$0.16$\,pc, $\lambda/W\approx1.4$; Section~\ref{sec:width}) show the sub-classical conclusion is robust to the convention, carried as the $\sim25\%$ width systematic (Section~\ref{sec:error_budget}).

\subsection{Width normalisation: $W_{\rm core}$ to $W_{\rm fil}$}
\label{sec:width_correction}

\begin{figure}
\centering
\includegraphics[width=0.99\columnwidth]{figures/fig_schematic.pdf}
\caption{Schematic of the length-scale definitions used throughout (Table~\ref{tab:symbols}). \emph{(a)} Transverse column-density profile, marking the initial half-width $W_{\rm core}=\sigma_0$, the evolved spine FWHM $W_{\rm form}$, and the observed beam-convolved Plummer FWHM $W_{\rm fil}=W_{\rm form}/T_1$ ($T_1\simeq0.61$; Eq.~\ref{eq:t1_apply}). \emph{(b)} Axial view contrasting the classical widely-spaced comb ($\lambda_{\rm m}\approx5$--$6\,W$) with the observed unevenly-distributed cores: a short intra-group adjacent spacing $s_{\rm ACS,med}$ within core-rich stretches separated by core-poor filament, whose filament-averaged line density $L/N$ is near-classical, while the simulated fragmentation scale is an upper bound $\lambda_{\rm turn}\lesssim1\,\lambda_{\rm J}$.}
\label{fig:schematic}
\end{figure}
Three widths enter the simulation--observation comparison and must be kept separate (Fig.~\ref{fig:schematic}; Table~\ref{tab:widthdefs}): the initial Gaussian half-width $W_{\rm core}=0.3\,\lambda_{\rm J}$; the formation FWHM $W_{\rm form}=0.683\,\lambda_{\rm J}$ of the evolved projected profile ($\simeq2.28\,W_{\rm core}$, slightly below the initial $2.355\,W_{\rm core}$ owing to contraction; resolution- and seed-converged to $<1\%$); and the observed HGBS width $W_{\rm fil}$, the beam-convolved Plummer FWHM \citep{Arzoumanian2019}, which we obtain by forward-modelling the simulated filaments through the full synthetic-HGBS pipeline (projection, beam convolution, Plummer fit). This gives $\eta_W=W_{\rm form}/W_{\rm fil}=0.59$--$0.78$ and
\begin{equation}
T_1 \equiv \frac{W_{\rm form}}{W_{\rm fil}} = 0.61 \quad (\text{forward-model band } 0.59\text{--}0.78),
\label{eq:t1_correction}
\end{equation}
so that $W_{\rm fil}=W_{\rm form}/T_1\approx1.12\,\lambda_{\rm J}$. We quote the forward-model best-fit $T_1=0.61$; because it lies near the lower edge of the band, a mid-band $T_1\approx0.68$ would give a slightly \emph{larger} (less sub-classical) $\lambda/W_{\rm fil}$, so this normalisation is deliberately not the headline; the primary comparison uses the width-convention-independent growth-rate upper bound (Section~\ref{sec:validation_endtoend}), and this $T_1$ conversion is retained only as a cross-check. The observable is $\lambda/W_{\rm fil}$; because the fragmentation spacing $\lambda$ is a physical length independent of the width definition, the conversion from the simulation-internal ratio $\lambda/W_{\rm core}$ (normalised by the initial $\sigma=0.3\,\lambda_{\rm J}$) is
\begin{equation}
\Big(\tfrac{\lambda}{W}\Big)_{\rm fil} = \frac{W_{\rm core}}{W_{\rm fil}}\Big(\tfrac{\lambda}{W}\Big)_{\rm core} \approx 0.27\,\Big(\tfrac{\lambda}{W}\Big)_{\rm core},
\label{eq:t1_apply}
\end{equation}
i.e.\ $\lambda/W_{\rm fil}=(\lambda/\lambda_{\rm J})/1.12$; the prefactor is $0.27=W_{\rm core}/W_{\rm fil}=0.30/1.12$ (equivalently $T_1$ times the $\sigma\!\to\!$FWHM factor, $0.61\times0.44=0.27$, using $T_1\,W_{\rm core}/W_{\rm form}$), and using $T_1$ alone would overstate $\lambda/W_{\rm fil}$ by $\approx2.3\times$. This prefactor is \emph{not} a fixed constant: it carries the full $\approx30\%$ $T_1$ band ($T_1=0.59$--$0.78$), so $0.27$ is the best-fit value of a $\approx0.26$--$0.34$ range. Because the observable is sensitive to how $W_{\rm fil}$ is defined, the primary comparison uses the width-convention-independent growth-rate wavelength (an upper limit; Section~\ref{sec:validation_endtoend}), and this normalisation is retained only as a cross-check.

\begin{table}
\caption{The three filament widths used in the simulation--observation comparison (all in Jeans lengths, $\lambda_{\rm J}=1$).}
\label{tab:widthdefs}
\footnotesize
\setlength{\tabcolsep}{3pt}
\begin{tabular}{p{1.5cm}p{1.4cm}p{4.0cm}}
\toprule
Symbol & Value ($\lambda_{\rm J}$) & Definition \\
\midrule
$W_{\rm core}$ & $0.30$ & initial Gaussian half-width $\sigma$ (input) \\
$W_{\rm form}$ & $0.683$ & Gaussian FWHM of evolved projected profile, supercritical ensemble ($\simeq2.28\,W_{\rm core}$) \\
$W_{\rm FWHM}$ & $0.58$ & Gaussian FWHM at the growth-rate epoch, near-critical runs (narrower, already contracting); normalises Table~\ref{tab:production} \\
$W_{\rm fil}$ & $\approx1.12$ & beam-convolved Plummer FWHM ($=W_{\rm form}/T_1$); the observable normalisation \\
\bottomrule
\end{tabular}
\end{table}

For the ambient-confined supercritical ensemble ($\lambda/W_{\rm core}=3.27$) this gives $\lambda/W_{\rm fil}\approx0.9$.

\subsection{Magnetic tension: a static-field estimate}
\label{sec:magnetic}

For a filament threaded by a magnetic field at an angle $\theta$ to the axis (Alfv\'en speed $v_A$), a schematic oblique dispersion relation, valid only for a static background field, is \citep{Nakamura1993}:
\begin{equation}
\omega^2 = c_s^2 k^2 + v_A^2 k^2 \sin^2\theta - 4\pi G\rho_0 F(kR)
\label{eq:dispersion}
\end{equation}
where $F(kR)$ encodes the cylindrical geometry and $\theta$ is the field--axis angle. This is a \emph{schematic static-field} relation that treats $\mathbf{B}$ as a fixed background and \emph{cannot be used quantitatively} for a magnetised cylindrical equilibrium; in particular it does not describe a current-carrying (helical) equilibrium. We use it only as an analytic guide, so this paper contains \emph{no quantitative model of magnetic tension}; it nonetheless fixes the \emph{sign} of each idealised case, the magnetic term $v_A^2k^2\sin^2\theta$ vanishing for a \emph{longitudinal} field (reducing to the field-free thermal relation) and being maximal for a \emph{perpendicular} one ($\theta=90^\circ$, which lengthens the axial mode). Consistent with this, our longitudinal runs keep the growth-rate wavelength at order $\lambda_{\rm J}$ across $\beta=0.3$--$1.0$ ($f=1.1$; the $\beta=1.0$ ladder gives $\lambda_{\rm turn}=1.14\,\lambda_{\rm J}$, Table~\ref{tab:production}) with no resolved $\beta$-trend, so the scale is set by the thermal Jeans length, not the field, and the shortening we measure is the dynamical/contraction shift, not magnetic tension. A genuinely shorter magnetic wavelength would require a force-balanced \emph{helical}/toroidal field, whose stability is a cylindrical magnetostatic-equilibrium problem that Eq.~(\ref{eq:dispersion}) cannot address. The full case-by-case discussion of the longitudinal, perpendicular and helical geometries, and the sense in which the simulations \emph{constrain only the specific initial field geometries simulated} rather than dispose of magnetic alternatives, is deferred to Section~\ref{sec:optional_tests}. The seeded helical run tested here is an \emph{out-of-equilibrium pinch}, not a force-balanced equilibrium, and does \emph{not} rule out the helical mechanism (Section~\ref{sec:optional_tests}).

\begin{table}
\caption{Schematic summary of the four magnetic-field configurations and their tested status.}
\label{tab:magcases}
\footnotesize
\setlength{\tabcolsep}{4pt}
\begin{tabular}{@{}p{2.9cm}lp{3.0cm}@{}}
\toprule
Configuration & Tested? & Effect on axial $\lambda$ \\
\midrule
Uniform longitudinal & tested & unchanged (term vanishes) \\
Static perpendicular (resolved) & tested & suppresses axial fragmentation (spindle) \\
Seeded/dynamic helical & tested & pinches $W$; does not shorten $\lambda$ \\
\midrule
\multicolumn{3}{l}{\emph{Not simulated in this work; listed for completeness}}\\
Force-balanced helical (Fiege--Pudritz) & not tested & could shorten by $\lesssim2\times$ (order-of-magnitude inference from Fiege--Pudritz, not a result of this work; Section~\ref{sec:optional_tests}) \\
\bottomrule
\end{tabular}
\end{table}

\subsection{Hierarchical fragmentation}

Filament fragmentation is hierarchical: filaments fragment into velocity-coherent fibres, and fibres into cores \citep{Hacar2013,Hacar2018,Yang2024}. Whether the fibre-to-core spacing is quasi-periodic at the classical scale is not established; the ALMA study of \citet{Yang2024} finds a clear filament--fibre--core hierarchy but only a weak quasi-periodic imprint at the core level. If HGBS filaments are bundles of fibres each fragmenting near the classical scale, measuring core spacings across the whole filament compresses the apparent spacing, in the observed direction. We test this quantitatively in Section~\ref{sec:fibre_projection}, where a \emph{toy geometrical model} of $N$-fibre bundles run through our ACS pipeline reproduces the observed $s_{\rm ACS,med}/W$; it emerges as a viable but not unique explanation (single-filament dynamical fragmentation is also consistent), and confirming it requires fibre-resolved core-spacing analysis beyond the single Orion~B region of \citet{Yang2024}.

\section{MHD simulations}
\label{sec:sims}

\smallskip\noindent\emph{The calculations below test whether sub-equilibrium fragmentation scales are dynamically possible. They are not intended to reproduce the observed core point process, and no numerical value quoted in them should be read as a prediction of an observed spacing.}

\noindent The calculations are exploratory. They are not a survey of filament fragmentation and they do not deliver a measured fragmentation wavelength: the growth-rate spectrum of a contracting filament has no resolved turnover, so every simulated number quoted here is an \emph{upper bound} $\lambda_{\rm turn}$ on where a turnover could lie. Their role is to provide an existence proof, that fragmentation from a contracting rather than an equilibrium state is not constrained to the classical wavelength, calibrated by the equilibrium-recovery control of Section~\ref{sec:validation_endtoend}. They should be read as a counterexample to a theoretical expectation, not as a prediction to be matched number-for-number against the observations.

The full campaign comprises $2251$ Athena++ runs across 18 grids, consuming of order $10^5$ CPU-hours on Google Cloud E2 infrastructure; controlled numerical experiments, not models of individual HGBS filaments. The large majority are \emph{exploratory} runs that map outcomes and timescales; the headline \emph{wavelength} result rests instead on the load-bearing \emph{production} runs (the 25-run near-critical growth-rate campaign, Table~\ref{tab:production}, and its 8-run $128^3$--$512^3$ convergence ladder; Table~\ref{tab:production}), so the $2251$-run total is \emph{not} the statistical weight behind the wavelength. Two results emerge (Section~\ref{sec:optional_tests}): the \emph{outcome} of a run (radial collapse or longitudinal beading) depends on regime and boundary conditions, so a single filament has no setup-independent fate, whereas the scale on which it fragments is bounded above at a sub-classical value ($\approx0.5$--$1\times$ the classical value; Section~\ref{sec:validation_endtoend}) for near-critical, longitudinal-field filaments (the perpendicular geometry is not numerically resolved). This section provides the evidence for both.

\subsection{Simulation method}
\label{sec:simmethod}

\smallskip\noindent\emph{The calculations below test whether sub-equilibrium fragmentation scales are dynamically possible. They are not intended to reproduce the observed core point process, and no number quoted in them is a prediction of an observed spacing.}

\smallskip We use Athena++ v21.0 \citep{Stone2020} with the \emph{isothermal} equation of state, $P=\rho c_s^2$. The energy equation is not evolved and no cooling prescription is applied; the evolved equations are continuity, momentum including the Lorentz force, and ideal induction, coupled to self-gravity through $\nabla^2\phi=4\pi G\rho$. The configuration uses the HLLD isothermal Riemann solver, piecewise-linear reconstruction in the primitive variables, the second-order van Leer integrator with CFL number $0.3$, constrained transport for the field, and the FFT Poisson solver with periodic gravitational boundaries along the axis and zero-gradient conditions transversely, matched to the hydrodynamic boundaries.

\smallskip\noindent\textbf{The initial condition.} The production runs use a Gaussian filament embedded in a uniform ambient medium,
\begin{equation}
\rho(r)=\rho_{\rm bg}+\rho_{\rm amp}\exp\!\left(-\frac{r^{2}}{2W_{\rm core}^{2}}\right),\qquad
\rho_{\rm amp}=\frac{f\,M_{\rm line,crit}}{2\pi W_{\rm core}^{2}},
\label{eq:ic_gauss}
\end{equation}
with $r^{2}=y^{2}+z^{2}$, $\rho_{\rm bg}=1$ in code units, $M_{\rm line,crit}=2c_s^{2}/G$ and $W_{\rm core}=0.3\,\lambda_{\rm J}$. The normalisation makes $\int_0^\infty(\rho-\rho_{\rm bg})\,2\pi r\,{\rm d}r=f\,M_{\rm line,crit}$ exactly, so the line mass converges without an imposed truncation and $f$ is set directly; the filament merges into the ambient medium, which supplies the confining pressure $P_{\rm bg}=\rho_{\rm bg}c_s^{2}$. The equilibrium-recovery control instead uses the Ostriker profile,
\begin{equation}
\rho(r)=\rho_{\rm bg}+\frac{f\rho_{c,\rm ost}}{\left[1+(r/W_{\rm core})^{2}\right]^{2}},\qquad
\rho_{c,\rm ost}=\frac{2c_s^{2}}{\pi G W_{\rm core}^{2}},
\label{eq:ic_ost}
\end{equation}
a Plummer profile of index $p=4$, whose line mass also converges. Neither uses a $p=2$ profile; the observed filaments have $p\simeq2$, which is why the width conversion of Section~\ref{sec:width_correction} is needed and why it is not invariant.

The remaining initial fields are $P(r)=\rho(r)c_s^{2}$, $v_r=v_\phi=0$, a uniform field of magnitude $B_0=c_s\sqrt{2\rho_{\rm bg}/\beta}$ at angle $\theta$ to the axis, and a longitudinal velocity seed of twelve Kolmogorov modes, $v_x=\sum_{n=1}^{12}A_n\cos(k_nx+\phi_n)$ with $k_n=2\pi n/L_x$, $A_n\propto k_n^{-11/6}$ and random phases. Domains are $L_x\times L_\perp^{2}$ with $L_x=16\,\lambda_{\rm J}$, periodic in $x$, and with the transverse boundary at $d\geq1\,\lambda_{\rm J}$ from the spine.

\smallskip\noindent\textbf{Parameters.} The filament is characterised by $f=M_{\rm line}/M_{\rm line,crit}$, the plasma $\beta=8\pi P/B^2$ and the seed-amplitude label $M_{\rm seed}=\delta v/c_s\times10^{4}$. At filament conditions ($n\sim10^4\,{\rm cm^{-3}}$, $T\sim10$\,K) $\beta=1$--$0.3$ corresponds to $B\sim15$--$30\,\mu$G \citep{Crutcher2012}. The mass-to-flux ratio scales as $\mu_\Phi/\mu_{\Phi,\rm crit}\propto f\sqrt{\beta}$ and exceeds unity throughout, so no configuration is magnetically subcritical. The load-bearing runs, listed in Table~\ref{tab:production}, are the near-critical growth-rate ladder and the direct supercritical ensemble; the wider exploratory grid is given in the deposited run manifest.

\smallskip\noindent\textbf{Weak seeds against transonic perturbations.} The production runs are seeded once, weakly and in the linear regime, whereas real filaments are turbulent at order-unity Mach number. A separate set of runs replaces the linear seed with an initial transonic velocity field at $\delta v/c_s=1$--$3$ drawn from the same spectrum. These shorten the time to beading by a factor $3$--$5$ but change the resulting wavelength by less than $5$ per cent (Pearson $r=-0.002$ between $\lambda/W$ and the perturbation Mach number) and do not introduce a maximum in $\Gamma(\lambda)$ where none was present, so the plateau is not an artefact of weak seeding. We retain the weak seed for the production runs because a growth-rate measurement requires a linear one: at transonic amplitude the perturbation is already non-linear when the measurement window opens. Continuously driven turbulence, in which energy is injected throughout the evolution, is not tested here.

\smallskip\noindent\textbf{Resolution.} The load-bearing runs satisfy the Truelove criterion at the spacing-setting epoch ($N_J\approx5$--$7$ cells per local Jeans length in both directions on the isotropic grids) but sit only $1.25$--$1.75\times$ above the $N_J=4$ minimum, which is one reason we quote a bound. Truelove stability and convergence are separate requirements: $N_J>4$ protects against artificial fragmentation but does not guarantee an accurate growth rate. The trustworthy band is therefore defined by the second criterion, $\Delta\Gamma/\Gamma<5$ per cent between successive resolutions, which here gives the same boundary.

\smallskip\noindent\textbf{Could numerical dissipation produce the missing maximum?} Grid-scale dissipation suppresses short-wavelength power progressively rather than at a sharp cutoff, and its signature is a short-$\lambda$ spectrum that \emph{rises} as the grid is refined. Three results argue against a dissipative origin for the plateau. The plateau edge does not move with resolution: over $128^3$--$512^3$ the fitted rate at the edge mode changes by $3.9$ and then $1.1$ per cent, while the spectrum below $\lambda\approx0.9\,\lambda_{\rm J}$ separates and rises, so the two regions behave in the opposite senses expected of a converged band and a dissipation-limited one. A numerical cutoff at fixed $\Delta x$ would move under refinement and would reveal a maximum hidden just below the band; it does not. And the axial mode set changes between resolutions, so a maximum pinned by quantisation would shift, and does not. The rows of Table~\ref{tab:production} that have been through this ladder are those at $f=1.1$ (three resolutions) and $f=1.5$ (two); the rest are single-resolution and quoted as bounds on that basis. What these tests exclude is dissipation removing a maximum lying \emph{within} the converged band, not a maximum below it.



\begin{table}
\caption{Load-bearing \emph{production} simulations underlying the fragmentation-wavelength result: the near-critical growth-rate ladder (campaign~15) and the direct supercritical ensemble (campaign~11). Shared parameters ($f$, $\mathcal{M}$, $\beta$, $\theta$, boundary conditions and collapse/bead time) are given in each block sub-header;  These production runs, not the wider exploratory grid, carry the quoted $\lambda_{\rm turn}$.}
\label{tab:production}
\footnotesize
\setlength{\tabcolsep}{4pt}
\begin{tabular}{ccccc}
\toprule
$f$ & RNG seed & resolution & $\lambda_{\rm turn}/\lambda_{\rm J}$ & $\lambda_{\rm turn}/W$ \\
\midrule
\multicolumn{5}{@{}p{8cm}}{\emph{Near-critical growth-rate ladder} (campaign~15): $\beta{=}1.0$, $\theta{=}0^\circ$, $\mathcal{M}{=}10^{-3}$ broadband seed, ambient BC (periodic/reflecting), bead $t_{\rm frag}{\approx}1.1$--$1.2\,t_{\rm J}$}\\
1.1 & 42  & $128^3$ & 1.14 & 1.95 \\
1.1 & 42  & $256^3$ & 1.14 & 1.99 \\
1.1 & 42  & $512^3$ & 1.14 & 2.00 \\
1.2 & 42  & $256^3$ & 1.00 & 1.81 \\
1.2 & 137 & $256^3$ & 0.89 & 1.60 \\
1.5 & 42  & $256^3$ & 0.89 & 1.47 \\
1.5 & 137 & $256^3$ & 0.89 & 1.47 \\
1.5 & 42  & $512^3$ & 0.89 & 1.47 \\
1.8 & 42  & $256^3$ & 0.80 & 1.19 \\
2.0 & 42  & $256^3$ & 0.80 & 1.20 \\
\midrule
\multicolumn{5}{@{}p{8cm}}{\emph{Direct supercritical ensemble} (campaign~11, 39 runs): $f{=}1.5$--$2.0$, $\beta{=}0.5$--$2.0$, $\theta{=}0^\circ$, seeds 42/137, periodic$+$reflecting ambient BC, bead $t_{\rm bead}{\approx}0.6$--$0.75\,t_{\rm J}$}\\
\multicolumn{2}{@{}c}{ensemble} & $256^3$ & $1.0$--$1.14$ & $0.89$--$1.01^{\dagger}$ \\
\bottomrule
\end{tabular}
\\[2pt]
\footnotesize $^{\dagger}$~In beam-convolved Plummer units, $\lambda/W_{\rm fil}$ (KS $p=0.17$ between periodic and reflecting BCs); the ladder's $\lambda_{\rm turn}/W$ column is in the Gaussian formation FWHM of Table~\ref{tab:production}.
\end{table}

\subsection{Why contraction can shift mode selection below the equilibrium scale}
\label{sec:whyupper}

Three ingredients make $\lambda_{\rm turn}$ a bound rather than a measurement.

\emph{(i) Contraction competes with mode growth.} The classical calculation assumes a static, marginally-stable cylinder, so every unstable mode has unlimited time and the winner is simply the one with the largest growth rate. Our production filaments begin with $f\gtrsim1$ and no supporting radial velocity, so the spine contracts while the axial modes are still growing. The relevant comparison is between the mode growth time $\tau_{\rm grow}=\Gamma(\lambda)^{-1}$ and the time on which the background changes, $\tau_{\rm contract}=|{\rm d}\ln\rho_c/{\rm d}t|^{-1}$; long-wavelength modes have the longest growth times and are the ones least likely to complete. The argument is set out in Appendix~\ref{sec:nullspec}.

\emph{Truelove-stable is not convergence-resolved.} These are separate requirements and we keep them separate. Satisfying $N_J>4$ protects against artificial fragmentation; it does not guarantee an accurate growth rate. The criterion we use to define the trustworthy band is the second one, $\Delta\Gamma/\Gamma<5$ per cent between successive resolutions, and the two happen to give the same boundary here. The bound $S_{\rm sim}^{\rm upper}\lesssim1.0$--$1.3$ should nonetheless be read as conditional on the resolution achieved: showing that a turnover has not appeared between two closely spaced marginal resolutions is not the same as excluding one below the numerical floor entirely.

\emph{(ii) There is therefore no resolved maximum.} Because the selection is set by the timescale competition rather than by the shape of $\Gamma(\lambda)$, the measured spectrum does not turn over: $\Gamma$ rises monotonically until it meets the grid scale, where the rise becomes numerical (Fig.~\ref{fig:growthrate}a). We quote the converged-band edge $\lambda_{\rm turn}$, an upper limit on where a physical turnover could lie, and never a peak.

\emph{(iii) The magnetic variables.} The field enters through $\beta$ and, equivalently, through $\mu_\Phi/\mu_{\Phi,\rm crit}\propto f\sqrt{\beta}$, which exceeds unity in every configuration run. In the uniform longitudinal-field configuration considered here a field does no work against axial modes, which compress gas along the field lines, and the simulations show no measurable dependence of the axial spectrum on $\beta$ over the tested range $0.3$--$2$. This is a statement about that configuration and that range, not a general result for longitudinally magnetised filaments. Its effect is radial: it resists transverse contraction and at $\beta\lesssim0.15$ halts it. A toroidal component could act on the axial mode, but only in radial force balance; seeded onto an already-contracting filament it merely pinches the spine, narrowing $W$ by $\approx25$ per cent and thereby \emph{raising} $\lambda/W$ without shortening $\lambda$.

\emph{Bridging to observable units.} Three widths are involved (Table~\ref{tab:widthdefs}): the initial Gaussian half-width $W_{\rm core}=0.3\,\lambda_{\rm J}$, the evolved formation FWHM $W_{\rm form}=0.683\,\lambda_{\rm J}$, and the beam-convolved Plummer FWHM $W_{\rm fil}$ that an observer fits. The bridge is one multiplicative factor,
\begin{equation}
\frac{\lambda_{\rm turn}}{W_{\rm fil}}=\eta_W\,\frac{\lambda_{\rm turn}}{W_{\rm form}},
\qquad \eta_W\equiv\frac{W_{\rm form}}{W_{\rm fil}}=T_1,
\label{eq:etaW}
\end{equation}
with $T_1=\eta_W=0.59$--$0.78$ measured by forward-modelling the simulation output through a synthetic \textit{Herschel} beam and the same Plummer fitting used on the data. It is Eq.~\ref{eq:etaW}, not the raw simulation number, that converts $\lambda_{\rm turn}/W_{\rm form}\approx2$ into $S_{\rm sim}^{\rm upper}\lesssim1.0$--$1.3$; the $\eta_W$ range is the dominant modelling term and widens that bound by about a quarter.



\subsection{Ambient confinement and the direct supercritical measurement}
\label{sec:modecompetition}

The outcome for a supercritical filament is set by whether it is confined by an ambient medium, as real HGBS filaments are ($d$ denotes the axis-to-boundary distance in Jeans lengths). The main 654-run grid has no ambient medium, so its Gaussian skirt feeds an accretion channel through the periodic boundary that suppresses beading and drives radial collapse; with ambient confinement (\texttt{filament\_ambient} generator) the same configurations fragment longitudinally, a 12-run reconciliation suite beading at physical, $\beta$-dependent collapse times ($0.55$--$0.70\,t_{\rm J}$) that converge across boundary conditions.

Ambient-confined supercritical filaments ($f=1.5$--$2.0$, $\theta=0^\circ$) fragment longitudinally: an ensemble of 39 runs (periodic and reflecting BCs, ambient confinement) develops 3--10 interior axial density maxima at a spacing of order the Jeans length. We take the wavelength from the linear growth-rate spectrum, whose low-$n$ modes are resolution-stable: no physical turnover is detected above $\approx1\,\lambda_{\rm J}$, so any turnover is constrained only to $\lambda_{\rm turn}\lesssim1.0$--$1.14\,\lambda_{\rm J}$ (an upper bound; no resolved peak), $\lesssim0.35$--$0.5$ times the classical value (not the median peak count, which is epoch- and resolution-dependent). The ambient medium confines the filament (central-to-ambient density ratio $\sim10^2$, consistent with HGBS crest-to-background contrasts; \citealt{Arzoumanian2019}), so it is representative rather than tuned. Across the 39 runs periodic and reflecting boundaries are statistically indistinguishable (KS $p=0.17$; medians $\lambda/W_{\rm fil}\approx0.89$ and $1.01$), and this ensemble median and the growth-rate result are the same physical wavelength $\lambda\approx1.1\,\lambda_{\rm J}$ in different normalisations. The ensemble beads early ($t_{\rm bead}\approx0.6$--$0.75\,t_{\rm J}$), before runaway and at moderate contrast, so with $\lambda\approx\lambda_{\rm J}$ resolved by ${\sim}32$ cells it satisfies Truelove at measurement (Appendix~\ref{sec:validation}).

Perpendicular-field filaments ($\theta=90^\circ$) do not universally spindle in the raw low-resolution grids: an 18-run ideal-MHD grid at $d=1.0$ shows $7/18$ bead (at $\beta=0.5$ and $2.0$) with an apparently short spacing ($\lambda/W_{\rm fil}\approx0.4$), but this conflicts with static linear theory (which predicts a \emph{longer} mode) and a transverse-refinement ladder to $N_J\approx170$ resolves the near-critical case into a monolithic radial spindle, identifying the short beading as a Truelove artefact (Section~\ref{sec:optional_tests}). An exploratory ambipolar-diffusion survey was also run, but it fragments only at $N_J\approx1$--$2$, below the Truelove limit, so it cannot support any conclusion and we exclude it from the argument here; it is described, with an explicit warning that it is numerically unresolved, in the deposited material.

\subsection{Testing the resolution and magnetic-field alternatives}
\label{sec:optional_tests}
Three targeted campaigns test the concerns that most affect the interpretation: whether the longitudinal wavelength is resolution-limited, whether the perpendicular geometry produces genuine short beading, and whether a helical field, the one configuration expected to shorten the axial mode, actually does.

\emph{Resolution stability of $\lambda_{\rm turn}$.} Two refinement directions behave oppositely. Under \emph{transverse} refinement (the isotropic $128^3$--$512^3$ ladder) the short-wavelength spectrum ($\lambda\lesssim0.9\,\lambda_{\rm J}$, $n\gtrsim9$) grows with resolution and is numerical (unresolved transverse sub-fragmentation; Fig.~\ref{fig:growthrate}a), whereas the physical band ($\lambda\gtrsim1\,\lambda_{\rm J}$, $n\leq8$) is converged. Under \emph{axial-only} refinement ($512\times128^2\to1024\times128^2$) the physical band reproduces to better than $0.5\%$ ($\Gamma$ at $n=6,7,8$, i.e.\ $\lambda=1.33,1.14,1.00\,\lambda_{\rm J}$, is $17.6,18.7,18.7$ at both resolutions, plateauing without turning over), \emph{and} the short-$\lambda$ rise is unchanged, confirming it is fixed by the axial Nyquist scale, not a physical mode migrating to shorter $\lambda$. The $\lambda_{\rm turn}\lesssim1\,\lambda_{\rm J}$ upper bound is therefore stable: axial refinement does not move the physical-band plateau, and the only shorter-$\lambda$ structure is the transverse-resolution artefact already identified.

\emph{The perpendicular geometry: a resolved spindle, not short beading.} On a transverse-refinement ladder ($128^3$--$512^3$) the near-critical \emph{ideal} perpendicular filament collapses monolithically in the radial direction at every resolution, the axial mode staying at the box scale and the axial perturbation at seed level throughout the well-resolved phase (at $512^3$ the spine stays resolved with $N_J\approx170$ cells at $C\approx3$, falling only to $N_J\approx40$ at $C\approx30$; still $\approx10\times$ above Truelove; Appendix~\ref{sec:validation}). Short-wavelength axial structure ($\lambda/W\approx0.4$) appears only once the spine drops below $N_J\approx1$; refinement delays its onset to deeper collapse, the signature of a grid artefact. The perpendicular field therefore \emph{suppresses} axial fragmentation in favour of radial collapse, as static linear theory predicts (a perpendicular field lengthens the axial mode), and the short low-resolution beading is numerical. This is for \emph{ideal} MHD; the ambipolar-perpendicular case (Section~\ref{sec:modecompetition}) remains separately unconverged.

\emph{A seeded helical field pinches the width but does not shorten $\lambda_{\rm turn}$; the force-balanced case remains untested.} A divergence-free helical-field scan ($B_\phi/B_z=0$--$2$; construction and scan detail in the deposited material) leaves the axial growth-rate plateau at $\lambda\approx1\,\lambda_{\rm J}$ for every winding and produces only a radial \emph{pinch} that narrows $W_{\rm FWHM}$ by $\approx25\%$ (raising $\lambda/W$ by reducing $W$, not by shortening $\lambda$; the deposited material), so, being seeded rather than force-balanced, it cannot test the confining Fiege--Pudritz sausage ($m=0$) mode \citep{Fiege2000}; the one geometry that could shorten $\lambda$, which therefore remains the single relevant magnetic configuration left unconstrained (Section~\ref{sec:optional_tests}).





\section{The simulated fragmentation scale}
\label{sec:validation_endtoend}

\subsection{What observable does a simulation predict?}
\label{sec:whatobservable}

A simulation that produces one bead per wavelength predicts a \emph{line density}, one core per $\lambda$ of filament. Its closest observational analogue is therefore $S_{\rm global}=(L/N)/W$ and not the conditional $S_{\rm local}$; the analogy holds only to the extent that the measured crest is a continuously unstable, velocity-coherent cylinder, which we have not established. The numerical proximity of the simulated bound ($\lesssim1.0$--$1.3$) to $S_{\rm local}\approx1.4$ is therefore not evidence of agreement, since the two are different quantities. Against $S_{\rm global}=3.4$--$5.8$ the simulated bound is in fact \emph{more} discrepant, implying that single-filament fragmentation over-produces cores per unit length unless most of the crest is core-poor, which is the occupancy result of Section~\ref{sec:intermittency}.

\subsection{The controlled experiment}
\label{sec:control}

The load-bearing numerical result is a control, not a campaign. We prepared a drag-relaxed, force-balanced Ostriker cylinder and a dynamically contracting filament and passed both through the \emph{identical} measurement pipeline. For the equilibrium the pipeline recovers the classical maximum at $22H$ ($\lambda_{\rm peak}/22H=1.07$--$1.10$ at two resolutions, a genuinely resolved maximum); for the contracting filament the same pipeline finds no maximum within the trustworthy band (Fig.~\ref{fig:eqrecovery}). The difference is a property of the prepared state, not of the estimator.

Because the argument rests on that distinction, the control is shown to be stationary rather than assumed to be. The cylinder is drag-relaxed to $t=6\,t_{\rm J}$ and then evolved drag-free. Over the last two Jeans times of the relaxation its gravitational binding energy changes by $0.6$ per cent and the mass-weighted radial velocity within $0.6\,\lambda_{\rm J}$ of the axis is $|v_r|_{\rm rms}\leq0.012\,c_s$; over the two Jeans times following removal of the drag, which is the interval in which growth rates are measured, the binding energy changes by $1.7$ per cent and $|v_r|_{\rm rms}\leq0.024\,c_s$. The residual drift is a slow expansion rather than a contraction, so it cannot mimic the contracting case, and if anything biases the control towards longer wavelengths, the conservative direction.

Wavelengths are measured from the growth-rate spectrum rather than by counting late-time peaks. Seeding twelve axial modes at $\delta v/c_s=10^{-3}$ and fitting $P(k,t)$ over its approximately exponential interval gives $\Gamma(\lambda)$ (Fig.~\ref{fig:growthrate}a), which is resolution-converged for $\lambda\gtrsim1\,\lambda_{\rm J}$ and rises with transverse resolution below $\lambda\approx0.9\,\lambda_{\rm J}$, where it is numerical. Within the converged band $\Gamma$ plateaus rather than turning over, so we quote only the upper bound $\lambda_{\rm turn}\lesssim1\,\lambda_{\rm J}$; forward-modelling through a synthetic beam and the same Plummer fitting used on the data (Eq.~\ref{eq:etaW}) gives the matched bound $S_{\rm sim}^{\rm upper}\lesssim1.0$--$1.3$.

Two caveats temper the causal attribution. The fitted rates ($\Gamma\approx18$--$26\,t_{\rm J}^{-1}$) are the accelerating growth of a collapsing background, not quasi-static eigenvalues, so we use only the mode \emph{ordering}; and a $p=4$ initial profile gives the same order-$\lambda_{\rm J}$ scale, so the result is not a profile-index artefact.

\begin{figure}
\includegraphics[width=0.99\columnwidth]{figures/fig_growthrate.pdf}
\caption{\emph{(a)} Growth-rate spectrum $\Gamma(\lambda)$ for a near-critical filament ($f=1.1$) at three resolutions. For $\lambda\gtrsim1\,\lambda_{\rm J}$ the curves overlap; at the grid scale (shaded) they separate and rise with resolution. $\Gamma$ plateaus rather than turning over, so $\lambda_{\rm turn}\lesssim1\,\lambda_{\rm J}$ is an \emph{upper bound}. \emph{(b)} $\lambda_{\rm turn}/W_{\rm FWHM}$ against line-mass fraction, in Gaussian FWHM units. The matched, beam-convolved Plummer bound is $S_{\rm sim}^{\rm upper}\lesssim1.0$--$1.3$ (grey band; the $\eta_W$ conversion is uncertain over $0.59$--$0.78$, widening it by about a quarter).}
\label{fig:growthrate}
\end{figure}

\begin{figure}
\centering
\includegraphics[width=0.99\columnwidth]{figures/fig_eqrecovery.png}
\caption{Equilibrium-recovery control. The identical growth-rate measurement pipeline is applied to a drag-relaxed, force-balanced pressure-confined Ostriker equilibrium (blue) and to a dynamically contracting filament (orange); the two differ in their prepared initial state, not in the measurement. \emph{Left:} onset fragmentation power spectrum; the equilibrium peaks at the classical $\lambda_{\rm m}\approx22H$ ($\approx5.4\,W_{\rm FWHM}$) while the contracting filament places its shortest converged scale at $\lesssim1.4\,W_{\rm FWHM}$ with no resolved maximum. \emph{Right:} $\Gamma(\lambda)$ against $\lambda/22H$; the equilibrium shows the classical peak-and-decline, the contracting filament rises monotonically. The pipeline recovers the classical wavelength when the physics is classical.}
\label{fig:eqrecovery}
\end{figure}

\subsection{The contraction ladder}
\label{sec:ladder}

The equilibrium-versus-contracting comparison establishes the sign of the effect but not its dependence on the state of the filament, so we ran a ladder. Seventeen otherwise identical filaments were initialised from the same relaxed Ostriker cylinder with a homologous inward velocity $v_r(r)=-A\,(c_s/R_0)\,r$, for $A$ from $0$ to $1.4$, at $256\times64^2$ with $\Delta x=0.0625\,\lambda_{\rm J}$ and $L_x=16\,\lambda_{\rm J}$, evolved with the same twelve-mode seed and the same growth-rate measurement, with the $A=0$ and $A=1$ endpoints repeated at $512\times128^2$.

The expectation to be tested is that mode selection follows the density at which the modes actually grow, $\lambda_{\rm select}\propto\rho_{\rm eff}^{-1/2}$, with $\rho_{\rm eff}$ the central density at the mid-point of the linear growth window. Figure~\ref{fig:ladder} shows why $A$ is not itself the controlling parameter. An imposed inward velocity does not produce steady contraction: the spine overshoots, reaching a transient central density up to $11$ times its initial value within $\approx0.3\,t_{\rm J}$, then rebounds and settles \emph{below} it, so that the most strongly kicked filaments are the least dense while their modes are growing. The relation between $A$ and $\rho_{\rm eff}$ is consequently non-monotonic, rising to $\rho_{\rm eff}\simeq3.5\,\rho_0$ near $A\approx0.2$ and falling to $0.35\,\rho_0$ at $A=1.4$. The ladder is therefore best read as a means of sampling $\rho_{\rm eff}$ over a factor of ten rather than as a sequence in contraction rate.

Across that range the selected scale follows the expected scaling. Three runs, at $A=0.15$, $0.2$ and $0.25$, show no resolved maximum within the seeded modes and enter the fit as upper limits; a maximum-likelihood regression that treats them as censored, and that includes the axial mode quantisation ($\lambda=L_x/n$) in the per-point uncertainty, gives
\begin{equation}
\lambda_{\rm select}\propto\rho_{\rm eff}^{-0.49\,\pm\,0.05},
\label{eq:ladderfit}
\end{equation}
with an intrinsic scatter of $0.09$ dex about the relation, against the predicted exponent of $-1/2$ (Fig.~\ref{fig:ladder}c). Discarding the censored points and fitting the remainder by least squares instead returns $-0.39$, which illustrates why the upper limits, all of which lie at the high-density end, must be included.

Evaluating the classical result at each run's own $\rho_{\rm eff}$ gives $\lambda_{\rm peak}/22H(\rho_{\rm eff})=1.09$ for the static case and $1.12$ for the weakest kick, both consistent with the classical value, and $0.79$--$1.04$ with a median of $0.86$ for the remaining fifteen. Contraction therefore shifts the selected mode modestly below the instantaneous classical value rather than dramatically so, and the shift is a consequence of the raised density rather than of the contraction itself. The $512\times128^2$ endpoints give $0.99$ and $0.84$ for $A=0$ and $A=1$, against $1.09$ and $0.94$ at the production resolution, so the ordering is unchanged under refinement.

\begin{figure*}
\includegraphics[width=0.98\textwidth]{figures/fig_ladder.pdf}
\caption{Contraction ladder, seventeen runs coloured by the imposed inward-velocity amplitude $A$. \emph{(a)} Central density against time, normalised to its initial value; the shaded band is the window over which the growth rates are fitted, and the circles mark the epoch at which $\rho_{\rm eff}$ is evaluated. The overshoot and rebound are visible directly: a strong kick raises the density transiently and leaves the filament less dense than it started while its modes are growing. \emph{(b)} Mass-weighted radial velocity within $0.6\,\lambda_{\rm J}$ of the axis, showing the same behaviour. \emph{(c)} Selected scale against $\rho_{\rm eff}$. Circles are runs with a resolved maximum of $\Gamma$; triangles are upper limits, where $\Gamma$ is still rising at the shortest well-resolved seeded mode. The solid line is the censored maximum-likelihood fit (Eq.~\ref{eq:ladderfit}) and the dashed line the $\rho_{\rm eff}^{-1/2}$ expectation.}
\label{fig:ladder}
\end{figure*}

Two things follow. The expected scaling is recovered over a factor of ten in density, which supports the interpretation that the wavelength is set by the instantaneous Jeans length at mode selection. But the controlling variable is that density and not the contraction rate: a filament can be given a large inward velocity and still select a \emph{long} wavelength if it has already rebounded by the time its modes grow. The mechanism should therefore be quoted as a dependence on $\rho_{\rm eff}$, and a real filament, which is fed continuously rather than kicked once, need not follow the same trajectory.

\subsection{What the simulations do and do not show}
\label{sec:simsummary}

The result is one-sided. These calculations rule out the assertion that a contracting filament \emph{must} fragment at the equilibrium-cylinder wavelength, and they show that mode selection follows the instantaneous central density. They do not determine a preferred wavelength, and they do not reproduce the observations: the simulated filaments produce a quasi-periodic comb of 3--10 near-uniform beads, whereas the observed cores reject both periodic and Poisson descriptions and occupy only $6$--$27$ per cent of the crest. A successful theory must predict a point process, not merely a wavelength, and no calculation presented here does so.

One further limitation is thermodynamic. Because $\Gamma(\lambda)$ does not turn over within the converged band, an isothermal calculation contains no small-scale cutoff and cannot predict a characteristic core mass or a CMF peak. That would require $\gamma_{\rm eff}>1$ at $n\gtrsim10^{5}$--$10^{6}\,{\rm cm^{-3}}$, where dust--gas coupling and line cooling can no longer hold $T$ constant \citep{Larson1985,Glover2012}; this is a change in thermal behaviour distinct from, and at much lower density than, the classical opacity limit \citep{Rees1976,LowLyndenBell1976}. Our isothermal suite deliberately excludes it.

\section{Discussion}

\noindent The observational material of Section~\ref{sec:results} onwards supports one composite statement: the cores in these four clouds do not form a quasi-uniform chain characterised by a single wavelength. They populate limited portions of the filament network, with short separations inside populated stretches and long core-poor intervals between them, and their gap distribution is inconsistent with both a periodic comb and a homogeneous Poisson process while showing no signature at the classical scale. The remainder of this section asks what could produce that pattern, and what the numerical experiments do and do not contribute to the answer.
\label{sec:discussion}

\subsection{Why the spacing matters: fragment masses and the CMF}
\label{sec:cmf}
The spacing matters because, with the line mass, it sets the characteristic fragment mass that seeds the core mass function (CMF) and ultimately the IMF: for a filament of line mass $M_{\rm line}$ fragmenting at spacing $\lambda$ the mean fragment mass is $M_{\rm frag}\sim\epsilon\,M_{\rm line}\,\lambda$, and a near-critical $M_{\rm line}\approx16\,M_\odot\,{\rm pc^{-1}}$ with $\lambda\approx0.1$--$0.2$\,pc gives $\sim0.5$--$1\,M_\odot$ for $\epsilon\sim0.3$, consistent with the HGBS CMF peak while the classical $5$--$6\times$FWHM would overshoot it. This is only an order-of-magnitude \emph{consistency check}: it is degenerate in $\epsilon M_{\rm line}$ (only the product enters) and yields a \emph{mean}, not a peaked distribution. An \emph{isothermal} EOS lacks a small-scale thermodynamic cutoff and therefore cannot predict a characteristic core mass or a unique CMF peak; only a non-isothermal EOS ($\gamma_{\rm eff}>1$ at $n\gtrsim10^{5}$--$10^{6}\,{\rm cm^{-3}}$) would convert the isothermal plateau into a physical peak (Section~\ref{sec:validation_endtoend}).

\subsection{Hierarchical-fibre projection}
\label{sec:fibre_projection}

A purely geometric route reaches a shorter local spacing without modified physics. HGBS filaments are bundles of $N$ velocity-coherent fibres \citep{Hacar2013,Hacar2018,Yang2024}; if each fragments near the equilibrium wavelength but spacing is measured \emph{across} the bundle, as any filament-level pipeline does, the interleaving compresses the apparent spacing. The premise, quasi-periodic fibre-to-core fragmentation, is a hypothesis rather than a measurement (\citealt{Yang2024} find only a weak imprint).

A toy model that reuses our exact measurement, populating a filament with $N$ cylinders of width $W_{\rm fib}$ each fragmenting at $\lambda_{\rm fib}\approx5$--$6\,W_{\rm fib}$ with random phase and $\sim30$ per cent lognormal scatter, gives in the interleaving limit
\begin{equation}
\frac{\lambda}{W_{\rm fil}}\approx\frac{(5\text{--}6)\,r}{N},\qquad r\equiv W_{\rm fib}/W_{\rm fil},
\label{eq:fourrn}
\end{equation}
verified across $r$ and $N$ in Fig.~\ref{fig:fibre_degen}. Either effect alone suffices: one fibre of half the filament width, or two full-width fibres, both give $\approx2.5$--$3.0$. The observation constrains only the combination $r/N$ ($\approx0.07$--$0.35$ across our conventions), so the relation cannot be inverted for the fibre multiplicity. Since a near-critical fibre would itself fragment below the equilibrium scale, the classical-fibre result is an upper bound and the fibre and single-filament pictures are not independent.

\begin{figure}
\includegraphics[width=0.48\textwidth]{figures/fig_fibre_degeneracy.pdf}
\caption{Hierarchical-fibre geometrical toy model, showing the filament-level $\lambda/W_{\rm fil}$ recovered by our ACS pipeline as a function of the fibre-to-filament width ratio $r=W_{\rm fib}/W_{\rm fil}$ and the number of fibres $N$; the model follows $\lambda/W_{\rm fil}\approx(5$--$6)r/N$ (Eq.~\ref{eq:fourrn}). White contours mark the observed fixed-width and region-specific values, the orange contour the FilFinder value, and the red band the observed fibre-width range ($r\approx0.3$--$0.5$). The observation constrains the combination $r/N$, not $N$ alone: the sub-classical value is reproduced by a single narrow fibre, by a few wider fibres, or by intermediate combinations.}
\label{fig:fibre_degen}
\end{figure}

Fibre projection is complementary to, not exclusive of, dynamical contraction, and the two are observationally degenerate in the filament-averaged spacing; breaking the degeneracy needs the tracer that resolves the fibres directly. A third route to the same reconciliation is post-formation evolution: cores accrete, migrate and interact over several free-fall times \citep{Kirk2016,Mattern2018}, so relaxation of an initially regular chain could erase it \citep{Clarke2020}. We do not simulate that stage.

\subsection{Field geometry and the untested magnetic configuration}
\label{sec:optional_tests}

Only a helical or toroidal field in radial force balance can shorten the axial mode. A uniform longitudinal field leaves it essentially unchanged in the configuration tested here; for a static perpendicular field the apparently short fragmentation is resolution-sensitive and is not established, the relevant runs having $N_J\approx1$--$2$, below the Truelove requirement, so they supply no constraint and are excluded from the argument; and a toroidal component seeded onto a contracting filament produces an out-of-equilibrium pinch rather than the confining equilibrium whose sausage ($m=0$) mode carries the shortened wavelength.

The force-balanced Fiege--Pudritz helical equilibrium, which is the one configuration that could genuinely shorten the axial mode, is not simulated here. Magnetic tension is therefore neither favoured nor disfavoured by these calculations, and the campaign places no constraint on this channel. Read off from the published \citet{Fiege2000} results rather than computed here, a strongly toroidally dominated equilibrium shortens the $m=0$ wavelength by up to a factor $\approx2$; our $\beta\approx0.3$--$2$ places these filaments at the weak end of that sequence. Settling the question needs a cylindrical MHD eigensolver: one would construct a Fiege--Pudritz self-similar equilibrium with poloidal $B_z(r)$ and toroidal $B_\phi(r)$ in radial force balance, confirm it is stationary by drag-relaxing and then evolving drag-free exactly as in our equilibrium control, and only then seed the $m=0$ mode.

\emph{Boundary conditions and non-ideal effects.} Periodic transverse boundaries supply an accretion channel that suppresses longitudinal fragmentation unless ambient confinement is imposed, which is why the production runs place the transverse boundary at $d\geq1\,\lambda_{\rm J}$ from the spine. At the densities probed here ($n_0\sim10^4$\,cm$^{-3}$) the ambipolar-diffusion time is long compared with the collapse time \citep{Tassis2004}, and non-ideal runs at $\eta_{\rm AD}=\{0.01,0.05\}$ still collapse radially, $12$--$15$ per cent earlier than the ideal control; ideal MHD is therefore an adequate approximation for the longitudinal stage.



\subsection{Limitations}
\label{sec:limitations}

Four limitations bound the conclusions.

\emph{Sample.} Four nearby, low-mass Gould Belt regions of broadly solar metallicity and similar radiation field are not a population, and whether the result persists in high-mass complexes, at lower metallicity or under stronger irradiation is untested. It is nevertheless worth asking how much of the cloud-to-cloud spread could be finite-sample noise. Drawing four samples of the observed sizes from the pooled gap distribution of all four clouds, the median range in $S_{\rm local}$ is $0.24$ and the $95$th percentile $0.48$, against an observed range of $1.31$ ($p=5\times10^{-5}$). The differences between the clouds in the conditional statistic are therefore real and not a small-$N$ artefact; what four clouds cannot support is the extrapolation of those differences to a population, or the identification of which cloud property drives them.

\emph{Method.} Two validation tests are outstanding. The $\lesssim5$ per cent bound on manual skeleton editing is a proxy, set by the junction-retention and association-radius sensitivities rather than by a blind multi-editor re-skeletonisation; and the small-separation deficit cannot be separated into instrumental and physical parts without a faint-segment injection--recovery experiment at the map level. All code and catalogues are deposited (Data Availability) so that both can be carried out.

\emph{Named follow-ups.} Three specific experiments would settle points left open here: projecting the three-dimensional dust reconstructions of \citet{Zucker2020} and \citet{Green2024} along the sightline through the Aquila/Serpens Rift, to replace the geometric blending model of Section~\ref{sec:methodology} by a count of physically distinct crests within one association length; injecting synthetic source pairs into the \textit{Herschel} maps and re-extracting them through the HGBS pipeline, to measure the close-pair completeness function and forward-model the gap distribution; and dense-gas position--position--velocity mapping to establish which skeleton corresponds to a velocity-coherent object.

\emph{Observations not available to us.} A per-segment line-mass split, which would test whether the most supercritical filaments show the shortest spacing, needs segment-level line masses at a fidelity the data do not support; resolved velocity dispersions per stretch, which would say why some eligible filament forms cores and some does not, do not exist for these regions.

\emph{Simulations.} These are idealised controlled experiments, not models of specific filaments. They are isothermal, use simplified initial profiles and a prescribed rather than driven turbulent spectrum, and the transverse boundaries are an idealisation of ambient confinement. The perpendicular-field and ambipolar parameter space is not numerically resolved at the resolutions we could afford ($N_J\approx1$--$2$, below the Truelove requirement) and therefore supplies no constraint on the fragmentation wavelength; those runs are excluded from the scientific argument and retained only in the deposited material.

\smallskip\noindent\emph{A recommendation for observers.} Because the skeleton systematic is the dominant limit on any comparison with theory, we recommend three things. First, that skeletons be extracted to a fixed column-density contrast relative to the local background rather than to an absolute level, so that the traced network does not depend on the distance or the mean column of the cloud. Second, that spacings be quoted for \emph{both} a persistent-crest and a medial-axis definition, so that the factor-of-two ontological span is stated rather than hidden inside a single number. Third, and most importantly, that any comparison with a one-dimensional cylinder model be based on velocity-coherent ridges identified in a dense-gas tracer rather than on a two-dimensional column-density mask: the classical calculation is an eigenmode of a single coherent object, and only kinematic data can establish that the traced segment is one. Where such data are unavailable, the persistent-crest network is the better of the two image-plane choices, because it follows the ridge rather than the extent of a mask, but the resulting comparison should be presented as a diagnostic and not as a test.



\subsection{Non-isothermal effects}
\label{sec:mt_variation}

The suite is isothermal, whereas HGBS dust-temperature maps vary by $\sim12$ per cent ($\sim10$--$13$\,K). Since $\lambda_{\rm J}\propto\sqrt{T}$ and $\mu_{\rm crit}\propto T$, temperature sets both the Jeans scale and the critical line mass, and an illustrative run with an imposed $\sim12$ per cent longitudinal gradient fragments differentially, the cool half beading while the warm half is suppressed, shifting the local wavelength by $\sim10$ per cent. This is a single run and indicates the direction and rough size of the effect only; a self-consistent treatment is deferred.



\section{Conclusions}

We have re-measured core spacing along filament crests in four HGBS clouds with \textit{Gaia} DR3 distances, and complemented the measurement with a controlled numerical experiment. Our conclusions are as follows.

\begin{enumerate}\setlength{\itemsep}{2pt}
\item \textbf{There is no single observational core spacing to compare with a one-wavelength theory.} The closest unconditional observational analogue of such a prediction is the filament length per core, $S_{\rm global}$, which is $3.4$, $3.4$, $4.5$ and $5.8$ in Orion~B, Aquila, Taurus and Perseus: essentially classical in Perseus, close to it in Taurus, and below the classical band by less than a factor of two in the other two. Measured only within core-bearing stretches, $S_{\rm local}$ is several times shorter.
\item \textbf{The difference between them is the result.} Cores occupy only $6$--$27$ per cent of the crest, forming tight groups separated by long bare stretches. Core formation along these filaments is therefore spatially intermittent, and what a fragmentation model must reproduce is a point process, not a spacing.
\item \textbf{The all-pairs median is the wrong estimator of local spacing.} It converges to $\approx0.293$ times the extent of whatever segment it is applied to. We use adjacent along-crest separations instead and retain the all-pairs statistic only for historical comparison.
\item \textbf{The magnitude of the local deficit is dominated by the choice of structural ontology, not by physics.} Changing from the persistent-crest network to a mask medial axis moves $S_{\rm local}$ by a factor of two and $S_{\rm global}$ by more. The \emph{sign} of the deficit is robust; the magnitude is not. Persistence thresholding within one ontology changes it by less than ten per cent (Section~\ref{sec:persistence}), so this is a question of definition rather than of parameter tuning, and it will be settled by velocity-resolved rather than image-plane data.
\item \textbf{The cores are neither periodic nor Poisson, but the smallest-separation deficit is not shown to be physical.} The adopted periodic-comb-plus-jitter model and the homogeneous Poisson model are both rejected in all four clouds, and the rejections survive a moving-block bootstrap. This is a rejection of those specific nulls, not of quasi-periodicity in general: a process with a spatially varying wavelength, missing fragments or subsequent migration would not be excluded by these tests. No classical-scale periodicity is detected in the pair-correlation function at the $95$ per cent level. The inferred exclusion scale, $0.011$--$0.044$\,pc, lies entirely within the \textit{Herschel} beam and \textsc{getsources} deblending limit, so it cannot be claimed as Jeans-scale suppression until map-level pair-injection tests are performed.
\item \textbf{The inhomogeneity is largely inherited from the filament.} Against a null conditioned on local column density we detect excess clumping in one cloud (Taurus) and fail to detect it in three. That null is conservative by construction, so the three non-detections are upper limits rather than evidence of absence, and with $N=4$ neither statement generalises to a population.
\item \textbf{The eligible-length diagnostic is exploratory.} Restricting to column-density-selected supercritical crest moves the two statistics together, but the criterion is partly endogenous and the correction is robust in only two clouds. None of the conclusions above depends on it.
\item \textbf{The simulations are a counterexample, not an explanation.} A contracting filament is not constrained to fragment at the equilibrium wavelength; the identical pipeline recovers the classical value for a relaxed equilibrium cylinder but finds no resolved maximum for a contracting one, so only an upper bound $S_{\rm sim}^{\rm upper}\lesssim1.0$--$1.3$ is available. A five-point ladder in contraction rate shows that mode selection tracks the density at which the modes grow, $\lambda_{\rm select}\propto\rho_{\rm eff}^{-0.55\pm0.12}$, close to the predicted $-1/2$. The simulations nevertheless produce a quasi-periodic comb and do not reproduce the observed intermittent point process, so they do not identify the origin of the observed spacing distribution.
\item \textbf{The physical cause is undetermined.} Dynamical contraction, hierarchical-fibre projection and helical magnetic tension all remain viable. The first is tested here, the second treated only geometrically, and the third not tested at all: no force-balanced helical equilibrium was simulated, so our campaign places no constraint whatever on the magnetic channel, which is mentioned for completeness rather than as a coequal candidate.
\end{enumerate}

The decisive next observations are velocity-resolved: dense-gas kinematics that identify which skeleton corresponds to a coherent physical object, and resolved velocity dispersions per stretch that would say why some supercritical filament forms cores and some does not.

\section*{Acknowledgements}

We thank the Herschel Gould Belt Survey team for the datasets. The Athena++ simulations were run on commercially procured cloud infrastructure (Google Cloud E2), funded from the author's institutional research allocation; the full set of $2251$ runs consumed of order $10^5$ CPU-hours.

\section*{Data availability}

The supporting material, comprising (i) the full run manifest (all 2251 runs with their $(f,\beta,\mathcal{M},\theta,\eta_{\rm AD})$ parameters and outcomes); (ii) the Athena++ problem-generator source (\texttt{filament\_ambient.cpp}, including the divergence-free helical construction) and normalisation scripts; (iii) HDF5 snapshots of the growth-rate convergence subset ($128^3$--$512^3$) and the equilibrium-recovery control (remaining snapshots on request); and (iv) the analysis pipeline (DisPerSE v0.9.24 parameters, the bridging/deblending and adjacent-core-spacing scripts, and the derived $\lambda/W$ catalogues); are archived under a CC-BY-4.0 licence with the citable Zenodo DOI \texttt{10.5281/zenodo.14872301}, versioned to the accepted release. The Zenodo record is the authoritative deposit and should be cited in preference to any working repository mirror. HGBS column-density and skeleton maps are available from the \textit{Herschel} Science Archive, and the \textit{Gaia} DR3 distances from \citet{Zhang2023}.

\bibliographystyle{mnras}
\appendix

\section{The limited HGBS sample}
\label{sec:limited_sample}
The four \emph{limited} regions fail at least one of the primary-selection criteria (Section~\ref{sec:distance_uncertainties}) and are excluded from the ACS measurement; they are retained only for the legacy pairwise-median comparison (Table~\ref{tab:literature}). the deposited material lists them together with the full 8-region core-count-weighted mean.



\section{The $L/3$ bias of the pairwise-median spacing}
\label{sec:pwmedian}
For $N$ points uniformly distributed on a filament of length $L$, a single pairwise distance $D=|X_i-X_j|$ has the triangular PDF $f(d)=2(L-d)/L^2$, with mean $\langle D\rangle=L/3$ and median $d^\star=L(1-1/\sqrt{2})\approx0.293\,L\approx0.88\,(L/3)$. The median of all $N(N-1)/2$ pairwise distances converges to $d^\star$ as $N\to\infty$, for a random or a perfectly periodic distribution alike, so the pairwise median measures the region extent, not the fragmentation wavelength. A Monte-Carlo over $N=5$--$500$ gives median$/L\approx0.293\,(1+0.43/N)$, i.e.\ $\le2\%$ inflation for $N\gtrsim20$, so the convergence is reached at the core counts of interest. For branched skeletons the statistic measures $\approx0.3\times$ the component's arclength extent, and for a realistic filament length distribution it tracks the rms-length$/3$ (full tables in the Zenodo archive). Computing the pairwise median over \emph{all} cores in each primary region (without per-filament segmentation) returns $\sim L_{\rm region}/3$ ($12.5$, $7.5$, $8.0$, $2.5$\,pc for Orion~B, Aquila, Perseus, Taurus), the region extent, confirming directly that the statistic measures segmentation scale rather than the fragmentation wavelength.

\section{Validation of the measurement pipeline}
\label{sec:validation}

The measurement chain was validated end to end before being applied. Rather than reproduce the individual tests here, we summarise them and note that the scripts and outputs are deposited with the code (Data Availability).

\emph{Closed-loop injection--recovery.} Synthetic core chains of known spacing were injected onto the real skeletons and recovered through the full pipeline; the recovered $\lambda/W$ matches the injected value to $0.4$ per cent over $\lambda/W=1.4$--$5.7$. A periodic injection over-corrects for a clustered process, which is why we quote raw values in the main text and carry the correction only as a bracket (Section~\ref{sec:injection_recovery}).

\emph{Eigenvalue check of the classical benchmark.} The dispersion relation of the isothermal equilibrium cylinder was solved directly and agrees with \citet{Nagasawa1987} and \citet{Fischera2012} to better than $1$ per cent, peaking at $\lambda_{\rm IM92}\approx2.3\,\lambda_{\rm J}$ with a cutoff at $\lambda_{\rm cr}\approx1.2\,\lambda_{\rm J}$. Because two normalisations are in use, we give the conversion explicitly. Throughout this paper $\lambda_{\rm J}\equiv c_s\sqrt{\pi/G\rho_0}$ is evaluated at the \emph{initial background} density $\rho_0$ (Table~\ref{tab:symbols}), whereas the scale height $H=c_s/\sqrt{4\pi G\rho_c}$ in Eq.~\ref{eq:lambdam} is evaluated at the cylinder's \emph{central} density $\rho_c$. For the equilibrium used here $\rho_c=2.24\,\rho_0$, so
\begin{equation}
\frac{22H}{\lambda_{\rm J}(\rho_0)}=\frac{22}{\sqrt{4\pi\times2.24}\times\sqrt{\pi}}=2.34,
\label{eq:hconv}
\end{equation}
which is the $2.3\,\lambda_{\rm J}$ quoted above. Had $\lambda_{\rm J}$ been evaluated at $\rho_c$ the same result would read $22H=3.50\,\lambda_{\rm J}(\rho_c)$, since $\lambda_{\rm J}=2\pi H$ at a common density. The independent solve therefore reproduces $22H$ exactly, and the two statements are the same number in different units. The measured contracting-filament spectrum instead rises monotonically into the grid-scale region, which is the basis for quoting an upper bound.

\emph{Width normalisation.} Forward-modelling the simulation output through a synthetic \textit{Herschel} beam and the same Plummer fitting used on the data gives $\eta_W=T_1=0.59$--$0.78$, with $W_{\rm form}=0.683\pm0.012\,\lambda_{\rm J}$ converged to $0.9$ per cent across resolutions. This is Eq.~\ref{eq:etaW}, the only route by which a simulated wavelength is compared with an observed spacing.

\emph{Resolution.} The load-bearing runs satisfy the Truelove criterion at the spacing-setting epoch ($N_J\approx5$--$7$), which is above the $N_J=4$ floor but not far above it, and this is one reason we quote a bound rather than a value. The growth-rate spectrum is resolution-converged for $\lambda\gtrsim1\,\lambda_{\rm J}$ at $128^3$--$512^3$ and separates only at the grid scale (Fig.~\ref{fig:growthrate}a).

\section{Robustness to the persistence threshold}
\label{sec:persistence}

The DisPerSE skeletons we use are the published HGBS products, extracted at a $3\sigma$ persistence cut. Since Section~\ref{sec:error_budget} identifies the skeleton algorithm as the largest single systematic, the persistence threshold of our primary ontology cannot be left untested. We cannot regenerate the published skeletons at a \emph{lower} threshold without the original DisPerSE run, so we approach the question from both sides.

\smallskip\noindent\emph{(i) Pruning the published skeleton upwards.} We decompose each skeleton into branches, terminating at junction pixels, and compute for each branch the persistence of its column-density maximum above the junction at which it attaches, in units of the local map noise $\sigma$ (estimated as $1.4826\times$MAD of the map minus a $5$-pixel median filter). Branches below a chosen threshold are removed and the procedure iterated to convergence. Raising the effective threshold from the published $3\sigma$ to $4\sigma$ removes $16$, $21$, $18$ and $15$ per cent of the skeleton length in Orion~B, Aquila, Perseus and Taurus, but changes the median adjacent spacing by only $-8.2$, $-2.7$, $+0.3$ and $-5.0$ per cent and the length per core by $-1.9$, $-3.2$, $-3.1$ and $-3.7$ per cent. Going further, to $6\sigma$ and $8\sigma$, changes nothing further at the per-cent level: the pruning saturates once the low-persistence spurs are gone. The occupancy fraction rises slightly, as expected when bare low-persistence branches are removed.

\smallskip\noindent\emph{(ii) An extraction with a tunable threshold.} Reaching the $2\sigma$ direction requires a crest extraction whose persistence threshold is a free parameter. We construct one directly from the column-density field: a descending-sweep merge tree pairs every maximum with the saddle at which it merges, maxima whose persistence falls below the chosen multiple of $\sigma$ are cancelled, and the surviving crests are traced by steepest-ascent walks from the retained saddles. This is not DisPerSE and does not reproduce the published skeleton pixel for pixel; it is used only to vary the threshold. Run at $2\sigma$, $3\sigma$ and $4\sigma$ on Orion~B it gives $S_{\rm local}=1.63$, $1.63$ and $1.63$ and $S_{\rm global}=2.53$, $2.50$ and $2.48$, with the associated core count falling from $635$ to $597$ and the occupancy fraction rising from $0.41$ to $0.42$. The local statistic is thus insensitive to the persistence threshold at the third significant figure over a factor of two in threshold, and the global statistic moves by two per cent.

\smallskip\noindent Both tests agree: the persistence threshold is not what makes the DisPerSE and FilFinder answers differ. Between $2\sigma$ and $8\sigma$ the local spacing moves by less than ten per cent, whereas changing the ontology from persistent crests to the medial axis of a mask moves it by a factor of two. The dominant systematic is therefore the \emph{definition} of a filament, not the tuning of a parameter within one definition, which is the conclusion of Section~\ref{sec:whatfilament}. The manual editing applied to the published skeletons cannot be varied in this way and remains an unquantified term.


\section{Sensitivity of the eligible-length diagnostic}
\label{sec:maskapp}

\begin{table}
\caption{Sensitivity of the core-masked eligibility diagnostic to the masking radius $r_{\rm mask}$. Cells give $S_{\rm elig}=(L_{\rm elig}/N_{\rm elig})/W$ with $N_{\rm elig}$ in brackets; $r_{\rm mask}=0$ is the uncorrected raw-map calculation and $0.06$\,pc is adopted. Orion~B ($1.56$--$2.04$) and Aquila ($3.14$--$3.21$) are comparatively stable; Perseus drifts and Taurus does not stabilise.}
\label{tab:maskrad}
\footnotesize
\setlength{\tabcolsep}{4pt}
\begin{tabular}{lcccc}
\toprule
$r_{\rm mask}$ (pc) & Orion~B & Aquila & Perseus & Taurus \\
\midrule
$0$ (raw) & 1.56 [417] & 3.14 [306] & 1.96 [199] & 0.81 [186] \\
$0.030$   & 1.89 [317] & 3.14 [306] & 3.68 [85]  & 1.80 [43] \\
$0.045$   & 2.04 [277] & 3.15 [305] & 3.91 [70]  & 1.61 [37] \\
$0.060$   & \textbf{2.02} [268] & \textbf{3.17} [303] & 4.13 [57] & --- [13] \\
$0.090$   & 1.98 [252] & 3.21 [296] & 4.45 [41]  & --- [1] \\
$0.120$   & 1.95 [238] & 3.20 [297] & 3.62 [37]  & --- [1] \\
\bottomrule
\end{tabular}
\end{table}

Table~\ref{tab:maskrad} gives the eligible line density of Section~\ref{sec:eligible} as a function of the radius within which column density is replaced around each associated core. Orion~B and Aquila are comparatively stable, varying by $\approx30$ and $\approx2$ per cent respectively over a factor of four in masking radius; Perseus drifts over $3.6$--$4.5$ while its eligible core count falls from $85$ to $37$; and Taurus does not stabilise at any radius. Only the first two clouds support the comparison.

\section{The inhomogeneous null, and a timescale interpretation of the contracting case}
\label{sec:nullspec}

\emph{The column-density-conditioned null.} The null used in Sections~\ref{sec:intermittency} and \ref{sec:clustering} asks whether cores are clumped \emph{beyond} the clumping already present in the filament they lie on. It is an inhomogeneous Poisson process on the one-dimensional skeleton with intensity
\begin{equation}
\lambda(x)\;\propto\;\left[\max\!\left(N_{\rm H_2}(x)-N_0,\,0\right)\right]^{\alpha},
\label{eq:nullintensity}
\end{equation}
where $x$ is arclength along the spine, $N_{\rm H_2}(x)$ the column density of the skeleton pixel at $x$, $N_0=10^{21}\,{\rm cm^{-2}}$ a floor that prevents low-column pixels from dominating, and $\alpha$ a single exponent fitted per sample. The normalisation is fixed by requiring $\int\lambda\,{\rm d}x=N_{\rm obs}$, so each realisation contains the observed number of cores. The exponent is fitted by maximum likelihood on the observed core positions, and, to avoid the circularity of fitting and testing on the same data, we refit it under leave-one-cloud-out cross-validation: for each cloud, $\alpha$ is estimated from the other three and applied to the held-out cloud. The fitted values are $\alpha=1.6$--$2.4$ with a cross-validated mean of $2.0$, and the conclusions are unchanged over $\alpha=1$--$3$.

For each cloud we draw $10^4$ realisations, compute the same Fano factor $F(1\,{\rm pc})$ and pair-correlation $g(s)$ estimators used on the data, and quote the empirical $p$-value $(1+\#\{F_{\rm null}\geq F_{\rm obs}\})/(1+N_{\rm MC})$ and the pointwise $2.5$th--$97.5$th percentile envelope. Two properties of this null should be kept in mind. It is conditional on the deposited skeleton, so it tests the distribution of cores along a fixed network and not the network itself. And it is \emph{conservative}: the column density entering Eq.~\ref{eq:nullintensity} is the observed map, which the cores themselves raise, so part of the signal being tested for is built into the null. That is the same endogeneity discussed in Section~\ref{sec:eligible}, and it biases the test against detecting excess clumping. Non-detections should therefore be read as upper limits on excess clumping rather than as evidence of its absence.

\emph{Why contraction can shift mode selection below the equilibrium scale.} The numerical result of Section~\ref{sec:sims} admits a simple timescale interpretation, which we give here because it is the reason we regard the counterexample as generic rather than accidental. It is an interpretation and not a derivation: the perturbations evolve on a time-dependent background and are not normal modes of a stationary equilibrium. In a filament of central density $\rho_c(t)$ the local Jeans length is $\lambda_{\rm J}(t)=c_s\sqrt{\pi/G\rho_c(t)}$, and an axial mode of fixed comoving wavelength grows at a rate $\Gamma(\lambda,t)$ that is positive only for $\lambda\gtrsim\lambda_{\rm J}(t)$. A mode is realised as a fragment only if it has time to grow before the background changes, that is if
\begin{equation}
\tau_{\rm grow}(\lambda)\;\equiv\;\Gamma(\lambda,t)^{-1}\;\lesssim\;\tau_{\rm contract}\;\equiv\;\left|\frac{\rm d\ln\rho_c}{{\rm d}t}\right|^{-1}.
\label{eq:freezeout}
\end{equation}
For an equilibrium cylinder $\tau_{\rm contract}\to\infty$ and every unstable mode has unlimited time, so the mode that dominates is the one with the largest $\Gamma$, the classical $\lambda_{\rm m}\approx22H$. For a contracting filament $\tau_{\rm contract}$ is of order the free-fall time and $\lambda_{\rm J}(t)$ falls as $\rho_c^{-1/2}$, so long-wavelength modes, whose growth times are the longest, do not complete before the state they were computed for has gone. The modes most likely to survive are those with $\tau_{\rm grow}\lesssim\tau_{\rm contract}$, and since $\Gamma$ increases towards shorter $\lambda$ over the converged band this selects the short end. The fragmentation scale is then set by the competition in Eq.~\ref{eq:freezeout} rather than by the maximum of $\Gamma$, which is why no turnover appears and why only an upper bound $\lambda_{\rm turn}$ can be quoted. The same argument predicts that the effect should weaken smoothly as the initial contraction rate is reduced towards zero, which is the ladder of Section~\ref{sec:ladder}.


\bibliography{references_complete}

\bsp
\label{lastpage}
\end{document}
